\section{Sampling order and the memory term of SAGA: i.i.d., fixed and reshuffled indices in a solvable model (hypothesis B4)}
\label{sec:B4_sampling_order}
% Sources of the numbers. The experimental numbers of the paragraph "Outcome of the experiments" are read in
% results/B4_summary.json (experiments/exp_B4_sampling_order.py) and results/B4b_summary.json
% (experiments/exp_B4b_refinements.py); the key is given in a comment before each item. The numbers printed elsewhere
% are values of closed forms of this note or of functions of vrlab/order_model.py and vrlab/bounds.py, named in the
% text; all of them are regenerated by experiments/exp_B4b_refinements.py under the entry "note" of
% results/B4b_summary.json (the key is given in a comment where the number is printed).
\providecommand{\Real}{\operatorname{Re}}

Remark~\ref{rem:B3-conclusion} of the note on Proposition~5.9 shows that no bound on $\E[V_t]$ can certify more than $\frac1n$ per iteration for SAGA with i.i.d.\ indices, i.e.\ one nat per pass for the squared error ($e^{-1}$ per pass): by \eqref{eq:B3-EM}, $\E[M_t]\ge(1-\frac1n)^tM_0$, because a memory slot that has never been refreshed keeps its initial error. The lower bound of the note on lower bounds for span methods with i.i.d.\ sampling has the same origin (Theorem~\ref{thm:instance}: a block that has never been sampled keeps all its energy, and the average unrevealed energy is $(1-\frac{1-q^2}n)^t$). The claim examined here (hypothesis (B4)) is that this ceiling is a property of sampling \emph{with replacement}, not of the memory table: an order that visits every index once per pass refreshes every slot once per pass. We make this exact on finite sums with identical Hessians, where SAGA reduces to a scalar linear recursion driven by the index sequence, and compare three orders at the same step: i.i.d.\ indices (R1), a fixed order (R2), and a fresh uniform permutation at every pass (random reshuffling, R3). Two sections then leave the model: one measures what a \emph{typical} run does under i.i.d.\ sampling, the other gives the constant that replaces the curvature when the Hessians differ and solves the opposite extreme, components on orthogonal directions. Throughout, exponents are \emph{per pass} ($n$ iterations) and for \emph{squared} errors, in nats: an exponent $r$ means a decay like $e^{-rt/n}$.

\begin{center}
\footnotesize
\begin{tabular}{@{}p{4.1cm}p{11.6cm}@{}}
\toprule
statement & status \\
\midrule
reduction to a scalar recursion (Lemma~\ref{lem:B4-reduction}) & proved, every $n$, every index sequence \\
(R1) i.i.d.: moment recursion, ceiling, plateau (Propositions~\ref{prop:B4-iid}, \ref{prop:B4-iid-limit}); rate at a fixed step (Proposition~\ref{prop:B4-iid-fixed}) & proved: exact at every $n$; limits $n\to\infty$ at fixed $c=\gamma n\lambda$ and at fixed $g=\gamma\lambda$ proved by perturbation of an eigenvalue \\
(R2) fixed order: recurrence, characteristic equation, growth rate (Lemmas~\ref{lem:B4-cyclic}, \ref{lem:B4-char}) & proved, every $n$ \\
(R2) limit $n\to\infty$ and critical step $c_\ast$ (Propositions~\ref{prop:B4-cyclic-limit}, \ref{prop:B4-critical}) & proved (Hurwitz's theorem and the argument principle); the finite-$n$ correction to $c_\ast$ and the best rate are formal or numerical (Remark~\ref{rem:B4-visible}) \\
(R3) reshuffling: first-moment pass matrix, its closed form and fluid limit, stability of the limiting first-moment map for every $c$ (Propositions~\ref{prop:B4-rr}, \ref{prop:B4-fluid}) & proved \\
(R3) exact second-moment pass matrix; at fixed $c$ the mean square decays at the rate of the squared mean (Propositions~\ref{prop:B4-rr2}, \ref{prop:B4-rr-ms}(i)) & proved: exact at every $n\ge2$; limit $n\to\infty$ proved \\
(R3) at a fixed step $g$: rate $\min\{2,\log\frac{2-g}g\}$, divergence of the mean square for $g>1$ (Proposition~\ref{prop:B4-rr-ms}(ii)) & proved (closed form of the weights of the fluctuation: the deviation is an autoregression of ratio $1-g$ driven by the memory, up to $O(n^{-2})$); checked on the exact matrix for $n$ up to $10^5$ \\
floors on the hard instance by order (Proposition~\ref{prop:B4-floors}) & proved (corollary of the note on lower bounds) \\
typical runs under i.i.d.\ sampling (Section~\ref{sec:B4-typical}) & existence of the typical exponent and the inequality ``typical $\ge$ expectation'' proved; its value is numerical, with a heuristic law \\
effective curvature for unequal Hessians; orthogonal components (Section~\ref{sec:B4-beyond}) & the identity behind the definition and the $2\times2$ maps of the orthogonal extreme are proved; the use of the effective curvature as a rate law is heuristic \\
logistic regression, hard instance (``Outcome'') & experimental \\
\bottomrule
\end{tabular}
\end{center}

\subsection{The model and the reduction to a scalar recursion}\label{sec:B4-model}

\begin{assumption}[identical Hessians]\label{ass:B4-model}
$f_i(\theta)=\frac12\theta^\top H\theta-b_i^\top\theta$ for $i=1,\dots,n$, with one symmetric positive definite $H\in\R^{d\times d}$, eigenvalues $0<\lambda_1\le\dots\le\lambda_d$ and orthonormal eigenvectors $w_1,\dots,w_d$, and arbitrary $b_i\in\R^d$. Then $F(\theta)=\frac12\theta^\top H\theta-\bar b^\top\theta$ with $\bar b=\frac1n\sum_ib_i$, $\etastar=H^{-1}\bar b$, and Assumption~\ref{ass:B3-setting} holds with $R^2=\lambda_d$, $\mu=\lambda_1$.
\end{assumption}

We run the SAGA iteration \eqref{eq:B3-saga} with a constant step $\gamma$ along an index sequence $(i(t))_{t\ge1}$ that may be random or deterministic, with the book's initialization $z_i^{(0)}=\fip(\theta_0)$. The four orders of \texttt{vrlab.algorithms.saga(..., sampling=)} are: \emph{i.i.d.} (uniform, independent draws); \emph{reshuffled} (passes of $n$ iterations, each using a fresh uniform permutation of $\{1,\dots,n\}$, independent of the past); \emph{shuffled once} (one uniform permutation, reused at every pass); \emph{cyclic} ($1,\dots,n$ at every pass). The last two are \emph{fixed orders}.

\begin{lemma}[reduction]\label{lem:B4-reduction}
Under Assumption~\ref{ass:B4-model}, let $\phi_j^{(t)}$ be the point at which the stored gradient $z_j^{(t)}$ was computed ($\phi_j^{(0)}=\theta_0$, $\phi^{(t)}_{i(t)}=\theta_{t-1}$, $\phi_j^{(t)}=\phi_j^{(t-1)}$ for $j\ne i(t)$), and write $e_t=\theta_t-\etastar$, $\psi^{(t)}_j=\phi^{(t)}_j-\etastar$, $\bar\psi_t=\frac1n\sum_j\psi_j^{(t)}$. Then, for every index sequence,
\begin{equation}\label{eq:B4-vector}
e_t=e_{t-1}-\gamma H\big(e_{t-1}-\psi^{(t-1)}_{i(t)}+\bar\psi_{t-1}\big),\qquad \psi^{(t)}_{i(t)}=e_{t-1},\qquad \psi_j^{(0)}=e_0 .
\end{equation}
Consequently each coordinate $\inner{w_l}{e_t}$, with the coordinates $\inner{w_l}{\psi^{(t)}_j}$ of the memory, follows the \emph{scalar recursion}
\begin{equation}\label{eq:B4-scalar}
e_t=e_{t-1}-g\big(e_{t-1}-\psi_{i(t)}+\bar\psi\big),\qquad\text{then}\quad \psi_{i(t)}\leftarrow e_{t-1},\qquad g=\gamma\lambda_l ,
\end{equation}
all modes being driven by the same index sequence, and $\norm{\theta_t-\etastar}^2=\sum_l\inner{w_l}{e_t}^2$.
\end{lemma}
\begin{proof}
$\fip(\theta)=H\theta-b_i$, so $z_j^{(t)}=H\phi_j^{(t)}-b_j$ by induction on $t$ and $\bar z_{t-1}=H\bar\phi_{t-1}-\bar b=H\bar\psi_{t-1}$ because $H\etastar=\bar b$. With $i=i(t)$, $\omega_t=H\theta_{t-1}-b_i-(H\phi_i^{(t-1)}-b_i)+H\bar\psi_{t-1}=H(e_{t-1}-\psi_i^{(t-1)}+\bar\psi_{t-1})$: the $b_i$ cancel. This is \eqref{eq:B4-vector}; taking the inner product with $w_l$ and using $Hw_l=\lambda_lw_l$ gives \eqref{eq:B4-scalar}.
\end{proof}

The $b_i$, hence the gradient noise at the optimum, have disappeared: the model isolates the memory term, and the only randomness left is the order of the indices. From now on we study \eqref{eq:B4-scalar} with the \emph{uniform memory} $\psi_j=e_0$ at time $0$, and measure the step in units of passes,
\[
c:=n\,g=\gamma\,n\,\lambda ,
\]
so that $n$ plain gradient steps would contract the mode by $(1-\frac cn)^n\approx e^{-c}$. For the textbook step $\gamma=\frac1{4R^2}$ the modes have $c_l=\frac{n\lambda_l}{4\lambda_d}\in[\frac n{4\kappa},\frac n4]$: the regime $\kappa\ll n$ of the discussion is the regime of large $c$.

\emph{Two scalings.} A mode is described by two numbers, $g=\gamma\lambda$ (the step, per iteration) and $c=ng$ (the step, per pass), and the limit $n\to\infty$ can be taken in two ways. At \emph{fixed $c$} the per-iteration step $g=\frac cn$ vanishes: this describes the modes whose $c$ is of order one, those that limit the rate when $\kappa$ is of the order of $n$ or larger. At \emph{fixed $g$} the per-pass step $c=ng$ diverges: this describes a stiff mode at a step that does not shrink with $n$, e.g.\ $g=\frac14$ for the stiffest mode at the textbook step, whatever $n$. Each statement below says which limit it takes; the exact finite-$n$ matrices cover both. The functions named below are those of \texttt{vrlab/order\_model.py}; each formula of this note was compared with them numerically, and the tests \texttt{test\_order\_model\_*} and \texttt{test\_saga\_identical\_hessian\_matches\_scalar\_model} of \texttt{tests/test\_vrlab.py} check the code against enumerations of all index sequences and against \texttt{saga} itself.

\subsection{(R1) I.i.d.\ indices: the ceiling of one nat per pass is attained in expectation}\label{sec:B4-iid}

\begin{proposition}[exact second moments, every $n$]\label{prop:B4-iid}
Let the indices be i.i.d.\ uniform. For the recursion \eqref{eq:B4-scalar} with any deterministic initial state, the moments $A_t=\E e_t^2$, $B_t=\E\frac1n\sum_j\psi_j^2$, $C_t=\E\bar\psi_t^2$, $P_t=\E e_t\bar\psi_t$ satisfy
\begin{equation}\label{eq:B4-moments}
\begin{aligned}
A_{t+1}&=(1-g)^2A_t+g^2(B_t-C_t), &
B_{t+1}&=\tfrac1nA_t+\big(1-\tfrac1n\big)B_t,\\
C_{t+1}&=\tfrac1{n^2}(A_t+B_t)+\big(1-\tfrac2n\big)C_t+\big(\tfrac2n-\tfrac2{n^2}\big)P_t, &
P_{t+1}&=\tfrac{1-g}nA_t-\tfrac gn(B_t-C_t)+(1-g)\big(1-\tfrac1n\big)P_t ,
\end{aligned}
\end{equation}
i.e.\ $(A,B,C,P)_{t+1}=T_n(g)\,(A,B,C,P)_t$ for a $4\times4$ matrix $T_n(g)$; the uniform memory is the initial vector $e_0^2\,(1,1,1,1)$. In particular $B_t\ge(1-\frac1n)^tB_0$ for every $t$.
\end{proposition}
\begin{proof}
Fix the state $(e,\psi)$ before a step and let $i$ be uniform; write $v=\frac1n\sum_j\psi_j^2-\bar\psi^2$. Then $\E_i[\psi_i-\bar\psi]=0$ and $\E_i(\psi_i-\bar\psi)^2=v$. Since $e^+=(1-g)e+g(\psi_i-\bar\psi)$, $\E_i(e^+)^2=(1-g)^2e^2+g^2v$. The memory update replaces $\psi_i^2$ by $e^2$ in the mean of squares: $\E_i\frac1n\sum_j(\psi_j^+)^2=\frac1ne^2+(1-\frac1n)\frac1n\sum_j\psi_j^2$. Next $\bar\psi^+=\bar\psi+\frac1n(e-\psi_i)$, so $\E_i(\bar\psi^+)^2=\bar\psi^2+\frac2n\bar\psi(e-\bar\psi)+\frac1{n^2}(e^2-2e\bar\psi+\frac1n\sum_j\psi_j^2)$. Finally $\E_i[e^+\bar\psi^+]=(1-g)e\bar\psi+\frac{1-g}ne(e-\bar\psi)+\frac gn\E_i[(\psi_i-\bar\psi)(e-\psi_i)]$ and the last expectation is $-v$. Taking expectations over the past gives \eqref{eq:B4-moments}; the last claim follows from $A_t\ge0$.
\end{proof}

\begin{proposition}[limit $n\to\infty$: ceiling and plateau]\label{prop:B4-iid-limit}
Fix $c>0$, let $g=\frac cn$ and $r^{\rm iid}_n(c):=-n\log\rho(T_n(\frac cn))$ ($\rho$ the spectral radius), the exponent of $\E e_t^2$ whenever the initial vector has a nonzero component on a dominant eigenvector (this holds for the uniform memory: by (b) for $c>\frac12$, and for $c<\frac12$ because the left eigenvector of $G$ for $-2c$ is $(1,0,0,0)$).
\begin{enumerate}[label=(\alph*)]
\item $n(T_n(\frac cn)-I)=G+\frac1nG'$ exactly, where, in the order $(A,B,C,P)$,
\[
G=\begin{pmatrix}-2c&0&0&0\\1&-1&0&0\\0&0&-2&2\\1&0&0&-c-1\end{pmatrix},\qquad
G'=\begin{pmatrix}c^2&c^2&-c^2&0\\0&0&0&0\\1&1&0&-2\\-c&-c&c&c\end{pmatrix}.
\]
The eigenvalues of $G$ are $-2c,-1,-2,-(c+1)$: the four modes of \eqref{eq:B4-moments} have exponents tending to $2c$, $1$, $2$ and $c+1$, and
\[
\lim_{n\to\infty}r^{\rm iid}_n(c)=\min\{2c,1\}=:r^{\rm iid}(c).
\]
\item Let $c>\frac12$. For $n$ large the spectral radius $\rho_n$ of $T_n(\frac cn)$ is a simple real eigenvalue, $\rho_n=1-\frac1n+\frac{c^2}{(2c-1)n^2}+O(n^{-3})$, so that $r^{\rm iid}_n(c)=1-\frac1n\big(\frac{c^2}{2c-1}-\frac12\big)+O(n^{-2})$, slightly \emph{below} one nat per pass; and for the uniform memory
\begin{equation}\label{eq:B4-plateau}
\E e_t^2=\alpha_n\,\rho_n^{\,t}\,e_0^2\,(1+o(1))\quad(t\to\infty),\qquad \lim_{n\to\infty}n\,\alpha_n=\frac{2c^3}{(2c-1)^2} .
\end{equation}
\end{enumerate}
\end{proposition}
\begin{proof}
(a) Multiply \eqref{eq:B4-moments} minus the identity by $n$ and insert $g=\frac cn$: for instance $n((1-g)^2-1)=-2c+\frac{c^2}n$ and $n((1-g)(1-\frac1n)-1)=-c-1+\frac cn$. In the order $(A,B,P,C)$ the matrix $G$ is lower triangular with diagonal $(-2c,-1,-c-1,-2)$. The eigenvalues of $T_n$ are $1+\frac{\nu}n$, $\nu$ an eigenvalue of $G+\frac1nG'$; eigenvalues depend continuously on the matrix, so these $\nu$ converge to the eigenvalues of $G$, and $-n\log\abs{1+\frac\nu n}\to-\Real\nu$. The smallest limit is $\min\{2c,1\}$ because $c+1>1$ and $2>1$.
(b) For $c\ne\frac12$ the eigenvalue $-1$ of $G$ is simple, with right eigenvector $v=(0,1,0,0)^\top$ and left eigenvector $w=(1,2c-1,0,0)$ (indeed $wG=(-2c,0,0,0)+(2c-1)(1,-1,0,0)=-w$). By the implicit function theorem applied to the characteristic polynomial, $G+\varepsilon G'$ has for small $\varepsilon$ a simple eigenvalue $\nu(\varepsilon)$, analytic and real, with $\nu(0)=-1$ and $\nu'(0)=\frac{wG'v}{wv}=\frac{c^2}{2c-1}$ ($G'v$ is the second column of $G'$), and an eigenvector $v(\varepsilon)\to v$ (the spectral projector of a simple eigenvalue is continuous), which we normalize by $v_B(\varepsilon)=1$. For $c>\frac12$ the three other eigenvalues of $G$ are smaller than $-1$, so for $n$ large $\rho_n=1+\frac1n\nu(\frac1n)$ is the unique eigenvalue of largest modulus of $T_n$; this gives the expansions of $\rho_n$ and of $-n\log\rho_n$. The component of the initial vector $\mathbf 1=(1,1,1,1)^\top$ on $v(\varepsilon)$ tends to $\frac{w\mathbf 1}{wv}=\frac{2c}{2c-1}$, and the first row of $(G+\varepsilon G')v(\varepsilon)=\nu(\varepsilon)v(\varepsilon)$ reads $(-2c+\varepsilon c^2)v_A+\varepsilon c^2(1-v_C)=\nu(\varepsilon)v_A$, whence $v_A(\varepsilon)=\varepsilon c^2(1-v_C(\varepsilon))/(2c+\nu(\varepsilon)-\varepsilon c^2)\sim\varepsilon\frac{c^2}{2c-1}$ since $v_C(\varepsilon)\to0$. Hence the coefficient $\alpha_n$ of $\rho_n^t$ in $A_t/e_0^2$ satisfies $n\alpha_n\to\frac{2c}{2c-1}\cdot\frac{c^2}{2c-1}$ (in particular $\alpha_n\ne0$ for $n$ large), and the other three modes are $o(\rho_n^t)$.
\end{proof}

\begin{remark}[reading (R1)]\label{rem:B4-iid-reading}
For $c<\frac12$ the expected squared error decays at the gradient rate $2c$; for $c\ge\frac12$ it cannot do better than one nat per pass, whatever the step. At leading order in $\frac1n$ the first row of $G+\frac1nG'$ is $A'=-2cA+\frac{c^2}n(B-C)$ in the time $u=\frac tn$: the error first decays like $e^{-2cu}$ and is then held, at the level $\frac1n$, by the dispersion $B-C$ of the memory, which decays like $\frac{2c}{2c-1}e^{-u}$ (solution of $B'=A-B$). The two regimes meet near $u_\times=\frac1{2c-1}\log\frac{n(2c-1)^2}{2c^3}$ (formal; for $n=1000$, $c=2$, the two-term approximation $e^{-2cu}+\frac{2c^3}{(2c-1)^2n}e^{-u}$ is within $1\%$ of the exact recursion at $u=8$ and $u=12$, and at $u=2\approx u_\times$ the exact value is $20\%$ below it). The plateau constant was checked on the exact recursion: at $n=20000$, $n\alpha_n=1.9996,\ 1.7776,\ 2.6121,\ 8.5240$ for $c=1,2,4,16$ (\texttt{iid\_slow\_amplitude}), against $2,\ 1.7778,\ 2.6122,\ 8.5245$.
% B4b_summary.json: note.R1_two_term_n1000_c2, note.R1_u_cross_n1000_c2, note.R1_n_times_amplitude_n20000 The inequality $B_t\ge(1-\frac1n)^tB_0$ of Proposition~\ref{prop:B4-iid} is the bound $\E[M_t]\ge(1-\frac1n)^tM_0$ of Remark~\ref{rem:B3-conclusion} in the model; what (R1) adds is that the \emph{error itself}, not only the Lyapunov function, inherits this rate, with an explicit constant.
\end{remark}

Proposition~\ref{prop:B4-iid-limit} is a limit at fixed $c$. At a step that does not shrink with $n$ the ceiling is lower.

\begin{proposition}[i.i.d.\ indices at a fixed step $g$]\label{prop:B4-iid-fixed}
Fix $g\in(0,2)$, $g\ne1$, and let $\varphi=\frac g{2-g}$. As $n\to\infty$ the eigenvalues of $T_n(g)$ are $(1-g)^2+O(\frac1n)$, $1-g+O(\frac1n)$, $1-\frac2n+O(n^{-2})$ and
\[
1-\frac{1-\varphi}n+O(n^{-2}),\qquad 1-\varphi=\frac{2(1-g)}{2-g}.
\]
Consequently, with $r^{\rm iid}_n$ as in Proposition~\ref{prop:B4-iid-limit}:
\begin{enumerate}[label=(\alph*)]
\item for $0<g<1$, $\lim_{n\to\infty}r^{\rm iid}_n(ng)=\frac{2(1-g)}{2-g}=:r^{\rm iid}_{\rm fix}(g)$ (\texttt{iid\_exponent\_fixed\_g}), which decreases from $1$ at $g=0$ to $0$ at $g=1$ and equals $\frac67$ at $g=\frac14$; for the uniform memory the coefficient of the dominant mode in $\E e_t^2/e_0^2$ tends to $\frac g2$;
\item for $1<g<2$, $T_n(g)$ has an eigenvalue larger than $1$ for $n$ large: the expected squared error diverges, although every plain gradient step contracts the mode ($\abs{1-g}<1$).
\end{enumerate}
\end{proposition}
\begin{proof}
By \eqref{eq:B4-moments}, $T_n(g)=T^{(0)}+\frac1nT^{(1)}+\frac1{n^2}T^{(2)}$ with, in the order $(A,B,C,P)$,
\[
T^{(0)}=\begin{pmatrix}(1-g)^2&g^2&-g^2&0\\0&1&0&0\\0&0&1&0\\0&0&0&1-g\end{pmatrix},\qquad
T^{(1)}=\begin{pmatrix}0&0&0&0\\1&-1&0&0\\0&0&-2&2\\1-g&-g&g&-(1-g)\end{pmatrix}.
\]
For $g\in(0,2)$ the eigenvalues $(1-g)^2$ and $1-g$ of $T^{(0)}$ differ from $1$, and the eigenvalue $1$ is semi-simple of multiplicity two: its right eigenvectors are $x_B=(\varphi,1,0,0)^\top$ and $x_C=(-\varphi,0,1,0)^\top$ (first row: $((1-g)^2-1)A+g^2(B-C)=0$, and $\frac{g^2}{1-(1-g)^2}=\varphi$), its left eigenvectors are the coordinate forms $e_B$, $e_C$ (the rows $B$ and $C$ of $T^{(0)}-I$ vanish), and $e_ix_j=\delta_{ij}$. By the first-order perturbation theory of a semi-simple eigenvalue \cite[Chapter~II, \S2.3]{kato1995perturbation}, the two eigenvalues of $T_n(g)$ near $1$ are $1+\frac\nu n+o(\frac1n)$, $\nu$ running over the eigenvalues of the reduced matrix
\[
\big(e_iT^{(1)}x_j\big)_{i,j\in\{B,C\}}=\begin{pmatrix}\varphi-1&-\varphi\\0&-2\end{pmatrix},
\]
i.e.\ $\nu=-(1-\varphi)$ and $\nu=-2$; these being distinct, the two branches are analytic in $\frac1n$, whence the $O(n^{-2})$. (a) For $0<g<1$, $0<1-\varphi<1<2$ and the two other eigenvalues have modulus at most $\abs{1-g}+O(\frac1n)<1$: the spectral radius is $1-\frac{1-\varphi}n+O(n^{-2})$ and $-n\log$ of it tends to $1-\varphi$. The left eigenvector of the reduced matrix for $\varphi-1$ is $(1,-\frac g2)$ (from $-\varphi-2\beta=(\varphi-1)\beta$), so the component of the initial vector $(1,1,1,1)$ on the dominant eigenvector tends to $1-\frac g2$, and its coordinate $A$ is $\varphi(1-\frac g2)=\frac g2$. (b) For $1<g<2$, $\varphi>1$ and the eigenvalue $1+\frac{\varphi-1}n+O(n^{-2})$ exceeds $1$.
\end{proof}

\noindent The mechanism is elementary. At fixed $g$ the error forgets its own past in $O(\frac1g)$ iterations, so it is \emph{slaved} to the memory: $e^+=(1-g)e+g(\psi_i-\bar\psi)$ has the stationary variance $\frac{g^2}{1-(1-g)^2}v=\varphi v$, where $v=B-C$ is the dispersion of the memory. A refreshed slot receives a value of variance $\varphi v$ in place of one of variance $v$, at the rate $\frac1n$ per iteration: $v$ decays at $1-\varphi$ nats per pass. We call
\[
\varphi(g)=\frac g{2-g}
\]
the \emph{fluctuation factor} (\texttt{fluctuation\_factor}). Proposition~\ref{prop:B4-iid-limit}(b) is consistent with this: its first-order correction $-\frac1n(\frac{c^2}{2c-1}-\frac12)$ is $-\frac g2+O(\frac1n)$ at large $c$, the expansion of $-\varphi(g)$, and its plateau $\frac{2c^3}{(2c-1)^2n}$ tends to $\frac g2$. Numerically (\texttt{iid\_exponent}): at $n=10^5$ the exact exponent is $0.857144$, $0.666666$ and $-0.222219$ for $g=\frac14,\frac12,1.1$, against $\frac67$, $\frac23$ and $-\frac29$.
% B4b_summary.json: note.R1_fixed_g; fixed_g.by_g.<g>.iid_exact

\subsection{(R2) A fixed order: instability beyond \texorpdfstring{$\gamma n\lambda=c_\ast\approx17.70$}{gamma n lambda = 17.70}}\label{sec:B4-cyclic}

\begin{lemma}[recurrence of a fixed order]\label{lem:B4-cyclic}
For the cyclic order and the uniform memory, the recursion \eqref{eq:B4-scalar} is, with the convention $e_k:=e_0$ for $k\le0$,
\begin{equation}\label{eq:B4-cyclic}
e_{t+1}=(1-g)\,e_t-\frac gn\sum_{s=1}^{n-1}e_{t-s}+g\Big(1-\frac1n\Big)e_{t-n}\qquad(t\ge0).
\end{equation}
Any fixed permutation repeated at every pass (in particular the shuffled-once order, whatever the permutation drawn) gives the same sequence $(e_t)$, hence, under Assumption~\ref{ass:B4-model}, the same iterates $\theta_t$.
\end{lemma}
\begin{proof}
At time $t\ge n$ the slot visited at step $t-s+1$ ($1\le s\le n$) holds $e_{t-s}$, so the memory is $\{e_{t-1},\dots,e_{t-n}\}$, its mean is $\frac1n\sum_{s=1}^ne_{t-s}$, and the slot to be visited next is the oldest one, $\psi_{i(t+1)}=e_{t-n}$. For $t<n$ the slots not yet visited hold $e_0$, which is $e_{t-s}$ for $s\ge t$ with the convention. Inserting into \eqref{eq:B4-scalar} gives \eqref{eq:B4-cyclic}. The recursion \eqref{eq:B4-scalar} and the uniform memory are invariant under a relabelling of the slots, which turns any fixed permutation into the cyclic order.
\end{proof}

Write $a_0=1-g$, $a_s=-\frac gn$ ($1\le s\le n-1$), $a_n=g(1-\frac1n)$ for the coefficients of \eqref{eq:B4-cyclic}, $p_n(x)=x^{n+1}-\sum_{s=0}^na_sx^{n-s}$ for its characteristic polynomial and $\rho_n(g)$ for the largest modulus of its roots (\texttt{cyclic\_radius}; in this section $\rho_n(g)$ always denotes this radius, not the spectral radius of $T_n$). Note $\sum_{s=0}^na_s=1-g$ and $\sum_{s=1}^na_s=0$.

\begin{lemma}[characteristic equation; every mode is excited]\label{lem:B4-char}
Let $n\ge2$, $g>0$ and $c=ng$.
\begin{enumerate}[label=(\roman*)]
\item With $y=n(x-1)$,
\begin{equation}\label{eq:B4-char-n}
n^2(x-1)\,p_n(x)=\Big(1+\frac yn\Big)^n\big(y^2+cy+c\big)-c\,(y+1)=:f_n(y),
\end{equation}
and $p_n(1)=g\ne0$: the roots of $p_n$ are the numbers $1+\frac yn$ with $y\ne0$ a zero of $f_n$.
\item For the uniform memory, $\limsup_{t\to\infty}\abs{e_t}^{1/t}=\rho_n(g)$. Hence $e_t\to0$ if and only if $\rho_n(g)<1$, and $(e_t)$ is unbounded if $\rho_n(g)>1$.
\end{enumerate}
\end{lemma}
\begin{proof}
(i) Since $(x-1)(x^{n-1}+\dots+x)=x^n-x$,
$(x-1)p_n(x)=x^{n+2}-(2-g)x^{n+1}+(1-g+\frac gn)x^n-gx+g(1-\frac1n)
=x^n\big[(x-1)^2+g(x-1)+\frac gn\big]-g(x-1)-\frac gn$; multiply by $n^2$ and use $n^2g(x-1)=cy$, $ng=c$. Moreover $p_n(1)=1-\sum_sa_s=g$.
(ii) Take $e_0=1$ and let $E(z)=\sum_{t\ge0}e_tz^t$. Multiplying \eqref{eq:B4-cyclic} by $z^{t+1}$, summing over $t\ge0$ and using $e_k=1$ for $k<0$ gives $E(z)-1=\sum_{s=0}^na_sz^{s+1}\big[E(z)+\sum_{k=1}^sz^{-k}\big]$, i.e.\ $E(z)D(z)=N(z)$ with $D(z)=1-\sum_{s=0}^na_sz^{s+1}=z^{n+1}p_n(1/z)$ and $N(z)=1+\sum_{s=1}^na_s(z+\dots+z^s)$. Then $(z-1)N(z)=z-1+\sum_{s=1}^na_sz^{s+1}-z\sum_{s=1}^na_s=z-1+(1-a_0z-D(z))=gz-D(z)$. A common zero $z_0$ of the polynomials $N$ and $D$ would satisfy $gz_0=0$, which is impossible since $D(0)=1$. So the rational function $E=N/D$ has a pole at every zero of $D$, i.e.\ at $1/x$ for every root $x$ of $p_n$ (these are nonzero, as $p_n(0)=-g(1-\frac1n)\ne0$), its radius of convergence is $1/\rho_n(g)$, and the first claim is the Cauchy--Hadamard formula. If $\rho_n<1$, $e_t\to0$ geometrically; if $\rho_n>1$ the sequence is unbounded; and if $\rho_n=1$ it does not tend to $0$, since the Taylor coefficients of a rational function with a pole on the unit circle and none inside do not tend to $0$ (by partial fractions they are a combination, with a nonzero coefficient for every pole, of the sequences $t^jx^t$ over the roots $x$ of $p_n$; the terms with $\abs x=1$ form a nonzero almost periodic sequence times a power of $t$, which does not tend to $0$, and the others do).
\end{proof}

\begin{proposition}[limit $n\to\infty$]\label{prop:B4-cyclic-limit}
For $c>0$ let
\begin{equation}\label{eq:B4-chi}
\chi_c(s):=s^2e^s+c\,(s+1)(e^s-1),\qquad \sigma(c):=\max\{\Real s:\ \chi_c(s)=0,\ s\ne0\}.
\end{equation}
\begin{enumerate}[label=(\roman*)]
\item The zeros of $\chi_c$ with $\Real s\ge-a$ satisfy $\abs s\le\max\{1,2c(1+e^a)\}$; $\chi_c$ has a real zero in $(-1,0)$; hence the maximum in \eqref{eq:B4-chi} exists and $\sigma(c)>-1$. The zero $s=0$ is simple.
\item $\lim_{n\to\infty}n\log\rho_n(\frac cn)=\sigma(c)$. Equivalently the squared-error exponent of a fixed order, $r^{\rm cyc}_n(c):=-2n\log\rho_n(\frac cn)$ (\texttt{cyclic\_exponent}), tends to $r^{\rm cyc}(c):=-2\sigma(c)$ (\texttt{cyclic\_exponent\_limit}).
\end{enumerate}
\end{proposition}
\begin{proof}
(i) If $\chi_c(s)=0$ then $e^s(s^2+cs+c)=c(s+1)$, so $\abs{s^2+cs+c}=c\abs{s+1}e^{-\Real s}\le ce^a(\abs s+1)$ and $\abs s^2\le c(1+e^a)(\abs s+1)$, which for $\abs s\ge1$ gives $\abs s\le2c(1+e^a)$. An entire function that is not identically zero has finitely many zeros in a compact set. On the real axis $\chi_c(-1)=e^{-1}>0$ and $\chi_c(s)=cs+O(s^2)<0$ for small $s<0$, so there is a zero in $(-1,0)$. Finally $\chi_c'(0)=c\ne0$.
(ii) By Lemma~\ref{lem:B4-char}(i) the roots of $p_n$ are the $x=1+\frac yn$ with $f_n(y)=0$, $y\ne0$, and $\abs x^n=\abs{1+\frac yn}^n$. If moreover $\abs x^n\ge e^{-1}$, the computation of (i) applied to \eqref{eq:B4-char-n} gives $\abs y\le Y:=\max\{1,2c(1+e)\}$, for every $n$. Since $(1+\frac yn)^n\to e^y$ uniformly on compact sets, $f_n\to\chi_c$ uniformly on compact sets. Choose $R>Y$ with no zero of $\chi_c$ on $\abs y=R$, let $s_0=0,s_1,\dots,s_K$ be the zeros of $\chi_c$ in $\abs y<R$, and let $\delta>0$ be so small that the discs $D(s_k,\delta)$ are disjoint and contained in $\abs y<R$. By Hurwitz's theorem there is $n_0$ such that for $n\ge n_0$, $f_n$ has in each $D(s_k,\delta)$ as many zeros as the multiplicity of $s_k$, and no other zero in $\abs y\le R$. As $s_0=0$ is a simple zero of $\chi_c$ and $f_n(0)=0$, the only zero of $f_n$ in $D(0,\delta)$ is $0$ itself, which is not a root of $p_n$. Hence, for $n\ge n_0$: every root $x$ of $p_n$ with $\abs x^n\ge e^{-1}$ has $y\in D(s_k,\delta)$ for some $k\ge1$, and every $D(s_k,\delta)$, $k\ge1$, contains a zero $y\ne0$ of $f_n$, i.e.\ a root of $p_n$. For $\abs y\le R$, $n\log\abs{1+\frac yn}=\Real y+O(R^2/n)$. By (i) with $a=1$, $\sigma(c)>-1$ is attained at some $s_k$ with $\abs{s_k}\le Y<R$: the root in its disc gives $n\log\rho_n\ge\sigma(c)-\delta+o(1)$; the roots with $\abs x^n\ge e^{-1}$ give at most $\max_{k\ge1}\Real s_k+\delta+o(1)=\sigma(c)+\delta+o(1)$, and the others have $n\log\abs x<-1<\sigma(c)$. As $\delta$ is arbitrary, the limit follows.
\end{proof}

The function $\chi_c$ is the characteristic function of the delay equation $e'(u)=-c\,e(u)+c\,e(u-1)-c\int_{u-1}^ue$, the formal limit of \eqref{eq:B4-cyclic} in the time $u=\frac tn$; Proposition~\ref{prop:B4-cyclic-limit} is what makes this limit rigorous for the growth rate.

\begin{proposition}[critical step of a fixed order]\label{prop:B4-critical}
Let $h(\omega):=\omega\sin\frac\omega2+\cos\frac\omega2$.
\begin{enumerate}[label=(\roman*)]
\item For $\omega\ne0$ real, $\chi_c(i\omega)=0$ if and only if $\omega^2=2c$ and $h(\omega)=0$, i.e.\ $\tan\frac\omega2=-\frac1\omega$. The positive zeros of $h$ are $\omega_1<\omega_2<\dots$ with exactly one $\omega_k$ in each interval $((2k-1)\pi,2k\pi)$ and none elsewhere. Set $c_k:=\frac12\omega_k^2$ and
\[
\omega_\ast:=\omega_1=5.9502,\qquad c_\ast:=c_1=\tfrac12\omega_\ast^2=17.702
\]
(\texttt{cyclic\_critical}; the next ones are $c_2=76.95$ and $c_3=175.65$).
\item At $c=c_k$ the zero $s=i\omega_k$ is simple, and the branch $s(c)$ through it satisfies
\begin{equation}\label{eq:B4-crossing}
\frac{ds}{dc}=\frac{2(1+s)}{2c^2+3c-2s},\qquad \Real\frac{ds}{dc}=\frac{2c\,(2c-1)}{(2c^2+3c)^2+8c}>0 .
\end{equation}
\item $\sigma(c)<0$ for $0<c<c_\ast$, $\sigma(c_\ast)=0$ and $\sigma(c)>0$ for $c>c_\ast$; more precisely $\chi_c$ has exactly $2k$ zeros in the open right half-plane for $c_k<c<c_{k+1}$ ($c_0:=0$).
\end{enumerate}
Consequently, by Proposition~\ref{prop:B4-cyclic-limit} and Lemma~\ref{lem:B4-char}(ii): for every $c<c_\ast$ a fixed order converges linearly on the mode for all $n$ large enough, and for every $c>c_\ast$ it diverges for all $n$ large enough. Under Assumption~\ref{ass:B4-model}, fixed-order SAGA with the initialization $z_i^{(0)}=\fip(\theta_0)$ (the uniform memory of Lemma~\ref{lem:B4-char}(ii)) therefore diverges as soon as one mode on which $\theta_0-\etastar$ has a component satisfies $\gamma n\lambda_l=c>c_\ast$ and $n\ge n_0(c)$.
\end{proposition}
\begin{proof}
(i) $e^{-i\omega/2}\chi_c(i\omega)=-\omega^2e^{i\omega/2}+2ic\,(1+i\omega)\sin\frac\omega2$, whose imaginary and real parts are $(2c-\omega^2)\sin\frac\omega2$ and $-\omega^2\cos\frac\omega2-2c\,\omega\sin\frac\omega2$. If $\sin\frac\omega2=0$ the real part is $\mp\omega^2\ne0$. Otherwise $\omega^2=2c$ and the real part is $-\omega^2h(\omega)$. With $\vartheta=\frac\omega2$, $h=2\vartheta\sin\vartheta+\cos\vartheta$ does not vanish where $\cos\vartheta=0$, is positive on $(0,\frac\pi2]$, and elsewhere vanishes iff $\ell(\vartheta):=\tan\vartheta+\frac1{2\vartheta}=0$; on each interval $((j-\frac12)\pi,(j+\frac12)\pi)$, $j\ge1$, $\ell'=\frac{2\vartheta^2-\cos^2\vartheta}{2\vartheta^2\cos^2\vartheta}>0$ and $\ell$ increases from $-\infty$ to $+\infty$, with $\ell(j\pi)>0$: exactly one zero, in $((j-\frac12)\pi,j\pi)$, i.e.\ $\omega=2\vartheta\in((2j-1)\pi,2j\pi)$.
(ii) At such a zero, $s^2=-2c$ gives $s^2+cs+c=c(s-1)$, so the equation $e^s(s^2+cs+c)=c(s+1)$ reads $e^s=\frac{s+1}{s-1}$ and $e^s-1=\frac2{s-1}$. Then $\partial_c\chi_c=(s+1)(e^s-1)=\frac{2(s+1)}{s-1}$ and $\partial_s\chi_c=e^s(s^2+cs+c+2s)+c(e^s-1)=\frac{(s+1)(c(s-1)+2s)+2c}{s-1}=\frac{-2c^2-3c+2s}{s-1}\ne0$. The implicit function theorem gives $\frac{ds}{dc}=-\partial_c\chi_c/\partial_s\chi_c$, and the real part of $2(1+i\omega)(2c^2+3c+2i\omega)$ is $2(2c^2+3c-2\omega^2)=2c(2c-1)$, positive since $c_k>\frac12$.
(iii) Let $N(c)$ be the number of zeros of $\chi_c$ in the open right half-plane, counted with multiplicity. \emph{Small steps.} Let $\chi_c(s)=0$ with $\Real s\ge0$, $s\ne0$, and write $e^s-1=sE(s)$, $E(s)=\int_0^1e^{\tau s}d\tau$, so $\abs{e^s-1}\le\abs se^{\Real s}$. From $s^2e^s=-c(s+1)(e^s-1)$ we get $\abs s\le c(\abs s+1)$, i.e.\ $\abs s\le r:=\frac c{1-c}$ for $c<1$. Dividing the equation by $s$: $se^s+c(s+1)E(s)=0$. For $s=x+iy$ with $x\ge0$ and $\abs s\le r<\frac\pi2$, $\Real(se^s)=e^x(x\cos y-y\sin y)\ge-e^rr^2$, and $\abs{(s+1)E(s)-1}\le\abs s\abs{E(s)}+\abs{E(s)-1}\le\frac32re^r$. Taking real parts, $0\ge c\,[1-e^rr(\frac32+\frac1{1-c})]$, which is false for $c\le\frac15$ ($r\le\frac14$, $e^{1/4}\cdot\frac14\cdot\frac{11}4<0.89$). So $N(c)=0$ for $c\le\frac15$. \emph{Continuity.} The function $\chi_c(s)/s$ is entire, analytic in $(c,s)$, equal to $c\ne0$ at $s=0$, and by (i) it has no zero on the imaginary axis unless $c\in\{c_k\}$; by Proposition~\ref{prop:B4-cyclic-limit}(i) with $a=0$ its zeros in the closed right half-plane lie in a disc whose radius is bounded on compact sets of $c$. The argument principle on the boundary of a large half-disc shows that $N$ is locally constant on $(0,\infty)\setminus\{c_k\}$. \emph{Crossings.} At $c=c_k$ the zeros on the axis are $\pm i\omega_k$, simple, and by \eqref{eq:B4-crossing} the two branches through them have $\Real s(c)$ of the sign of $c-c_k$ near $c_k$, while the other zeros of the closed half-plane stay in the open one: $N$ increases by $2$ at each $c_k$. Hence $N=2k$ on $(c_k,c_{k+1})$; $N=0$ and the absence of zeros on the axis give $\sigma(c)<0$ on $(0,c_\ast)$, and $N\ge2$ gives $\sigma(c)>0$ on $(c_\ast,\infty)\setminus\{c_k\}$, the values $\sigma(c_k)$ being $0$ for $k=1$ and positive for $k\ge2$.
\end{proof}

\begin{remark}[how visible is the instability; finite $n$; best rate]\label{rem:B4-visible}
\emph{(a) The instability is weak (proved to first order, then numerical).} By \eqref{eq:B4-crossing}, $\sigma(c)=0.002635\,(c-c_\ast)+O((c-c_\ast)^2)$: just above the threshold the squared error grows by $2\sigma\approx0.0053\,(c-c_\ast)$ nats per pass. Numerically (\texttt{cyclic\_exponent\_limit}; the rightmost zero is found by Newton's method from a grid of starting points, not by a certified count) $r^{\rm cyc}(c)=-0.0015,\ -0.0097,\ -0.0244,\ -0.0240,\ -0.0157,\ -0.0046$ at $c=18,20,30,50,100,400$, and over $c\in[17.8,400]$ the fastest growth is $2\sigma\approx0.026$ nats per pass, near $c=37$; for large $c$ the zeros near $2\pi ik$ have $\Real s\approx\frac{4\pi^2k^2}{c(1+4\pi^2k^2)}$ (formal expansion; $\Real s=9.56\cdot10^{-4}$ against $9.75\cdot10^{-4}$ for $k=1$, $c=10^3$). In the limit model a fixed order beyond $c_\ast$ therefore multiplies the squared error by at most about $e^{0.026}$ per pass (at finite $n$ a little more: $0.035$ nats per pass at $n=100$ and $120$ for $g=\frac14$, see (b)): over the tens of passes of an experiment this is a \emph{stall} of the stiff modes rather than an explosion, which is how it should be looked for in practice.
% B4b_summary.json: note.R2_slope_at_c_star, note.R2_exponent_limit, note.R2_fastest_growth_on_17.8_400, note.R2_large_c_root_near_2pi_c1000
\emph{(b) Finite $n$ (formal, then numerical).} On the unit circle $x=e^{s/n}$ with $s$ imaginary and $y=s+\frac{s^2}{2n}+O(n^{-2})$; inserting into \eqref{eq:B4-char-n} and using the identities of the proof of (ii), the zero at $c=c_\ast$ is shifted to first order by $\delta s$ with $\Real\delta s=\frac{2c^2(2c^2+3c+4)}{n((2c^2+3c)^2+8c)}>0$, so the critical step at finite $n$ is
$c_\ast(n)=c_\ast-\frac Kn+O(n^{-2})$ with $K=\frac{c_\ast(2c_\ast^2+3c_\ast+4)}{2c_\ast-1}\approx351.9$ (the remainder is not controlled here). The companion matrix gives $c_\ast(n)=12.716,\ 14.784,\ 16.107,\ 16.866,\ 17.273,\ 17.358$ for $n=50,100,200,400,800,1000$, i.e.\ $n(c_\ast-c_\ast(n))=249,\ 292,\ 319,\ 335,\ 343,\ 345$: the threshold is approached from below, slowly. In particular, for the textbook step the stiffest mode has $g=\frac14$, and $\rho_n(\frac14)>1$ for every $n$ from $52$ to $120$ and for $n=200$ and $1000$ (\texttt{cyclic\_radius}, checked for each of these $n$; $\rho_n(\frac14)<1$ for $2\le n\le51$): \emph{for these $n$, on an identical-Hessian sum, cyclic or shuffled-once SAGA with $\gamma=\frac1{4R^2}$ does not converge} unless $\theta_0-\etastar$ has no component on the top eigenspace of $H$ (Lemma~\ref{lem:B4-char}(ii)), whereas the i.i.d.\ and reshuffled runs do: on \texttt{identical\_hessian\_problem(100, [0.02, 0.3, 1], seed)}, $\theta_0=0$, the squared error of the two fixed orders (identical iterates up to rounding) is $4.0$, $122$ and $10.7$ at the start for the seeds $0,1,2$, $2.61$, $0.049$ and $0.97$ after $100$ passes and $219$, $4.15$ and $81.2$ after $300$: after an initial drop the stiff mode grows by $e^{0.035}$ per pass, while the two other orders are below $10^{-25}$ after $100$ passes.
% B4b_summary.json: note.R2_K, note.R2_threshold_finite_n, note.R2_textbook_step_g_quarter, note.R2_example_n100_eigs_0.02_0.3_1
\emph{(c) Best rate (numerical).} $r^{\rm cyc}(c)$ is maximal at $c=2.148$, where the real zero and a complex pair have the same real part, with value $1.526$ (\texttt{cyclic\_best}); it is $0.756$, $1.149$, $1.497$, $0.702$, $0.194$, $0.011$ at $c=\frac12,1,2,4,8,16$.
\emph{(d) Incremental aggregated gradient (exact; the limit is a classical delay equation).} The same computation for the direction $\bar z$ alone (IAG \cite{blatt2007convergent,gurbuzbalaban2017convergence}; in the model $e_{t+1}=e_t-\frac gn\sum_{s=0}^{n-1}e_{t-s}$) is explicit at every $n$: the characteristic polynomial is $x^n-x^{n-1}+\frac gn(x^{n-1}+\dots+1)$, and, after multiplication by $x-1$, a root $x=e^{i\vartheta}\ne1$ on the unit circle requires $e^{in\vartheta}=-1$ and $\frac gn=2\sin^2\frac\vartheta2$, first at $\vartheta=\frac\pi n$: the critical step is
\[
c^{\rm IAG}_\ast(n)=2n^2\sin^2\frac\pi{2n}\ \xrightarrow[n\to\infty]{}\ \frac{\pi^2}2\approx4.93
\]
(\texttt{iag\_critical}; $4.93440$ at $n=100$, equal to the value at which the companion radius \texttt{iag\_radius} crosses $1$, and $4.93480$ at $n=1600$). The limit is the stability boundary of the delay equation $e'(u)=-c\int_{u-1}^ue$, with characteristic equation $s^2+c(1-e^{-s})=0$, which is classical: its zero solution is asymptotically stable if and only if $0<c<\frac{\pi^2}2$ \cite[Remark~1.1 and the references given there]{matsunaga2011asymptotic}. For the method itself \cite{schmidt2017minimizing} report that IAG with a step $\frac1{16L}$ may diverge and that a step of about $\frac1{nL}$ is needed, the exact value being ``not precisely characterized''; in the model it is $\gamma n\lambda<c^{\rm IAG}_\ast(n)$. The correction $f'_i(\theta)-z_i$ of SAGA moves the threshold of a fixed order from $4.93$ to $17.70$, but does not remove it.
% B4b_summary.json: note.IAG_threshold
\end{remark}

\subsection{(R3) Random reshuffling: exact first moments, two nats per pass}\label{sec:B4-rr}

\begin{proposition}[exact first moments over a reshuffled pass, every $n$]\label{prop:B4-rr}
Consider one pass of \eqref{eq:B4-scalar} along a uniform random permutation that is independent of the state $(e,\psi)$ at the start of the pass; let $m=\bar\psi$ be the mean of the memory at that time, $e'$ the error at the end of the pass and $m'$ the mean of the memory at the end of the pass. Then
\[
\E\big[(e',m')^\top\,\big|\,e,\psi\big]=M_n(g)\,(e,m)^\top ,
\]
where the $2\times2$ matrix $M_n(g)$ (\texttt{rr\_pass\_matrix}) is defined by the recursion $\bar e_0=e$, $S_0=0$,
\begin{equation}\label{eq:B4-rr-rec}
\bar e_{k+1}=(1-g)\,\bar e_k-\frac gn\,S_k+\frac{g\,k}n\,m,\qquad S_{k+1}=S_k+\bar e_k\quad(0\le k<n),\qquad (e',m')\mapsto\big(\bar e_n,\tfrac1nS_n\big).
\end{equation}
In closed form, with $\varpi=1-g+\frac gn$ and the sequence $U_0=0$, $U_1=1$, $U_{k+1}=(2-g)U_k-\varpi\,U_{k-1}$ (i.e.\ $U_k=\frac{r_+^k-r_-^k}{r_+-r_-}$, $r_\pm=1-\frac g2\pm\sqrt{\frac{g^2}4-\frac gn}$, when $g\ne\frac4n$; the recurrence defines $U_k$ in every case),
\begin{equation}\label{eq:B4-rr-closed}
M_n(g)=\begin{pmatrix}U_{n+1}-U_n & 1+(n-1)(U_{n+1}-U_n)-n\varpi\,(U_n-U_{n-1})\\[2pt] \frac1nU_n & \frac1n\big((n-1)U_n-n\varpi\,U_{n-1}\big)\end{pmatrix},
\end{equation}
and $\det M_n(g)=\varpi^n-\frac1nU_n$.
Consequently, for reshuffled SAGA started from the uniform memory, $\E[e_{pn}]=\big[M_n(g)^p(1,1)^\top\big]_1e_0$ exactly for every number $p$ of passes, and
\begin{equation}\label{eq:B4-rr-jensen}
\E\,e_{pn}^2\ \ge\ \big(\big[M_n(g)^p(1,1)^\top\big]_1\big)^2e_0^2 .
\end{equation}
\end{proposition}
\begin{proof}
Let $e_k$ be the error after $k$ steps of the pass, $\pi(k+1)$ the slot visited at step $k+1$ and $S_k=\sum_{j<k}e_j$. A slot is visited once per pass, so at step $k+1$ the slot $\pi(k+1)$ still holds its value $\psi_{\pi(k+1)}$ from the start of the pass, the $k$ slots already visited hold $e_0,\dots,e_{k-1}$, and
$e_{k+1}=(1-g)e_k+g\,\psi_{\pi(k+1)}-\frac gn\big(S_k+\sum_{j\notin\{\pi(1),\dots,\pi(k)\}}\psi_j\big)$.
Conditionally on $(e,\psi)$, $\pi(k+1)$ is uniform on the slots, so $\E\psi_{\pi(k+1)}=m$, and each slot is unvisited after $k$ steps with probability $\frac{n-k}n$, so the last sum has expectation $(n-k)m$. Only linearity of the expectation is used (the $e_k$ depend on the permutation): with $\bar e_k=\E[e_k\mid e,\psi]$ and $S_k$ now denoting $\E[S_k\mid e,\psi]$, $\bar e_{k+1}=(1-g)\bar e_k+gm-\frac gn(S_k+(n-k)m)$, which is \eqref{eq:B4-rr-rec}. At the end of the pass the memory is $\{e_0,\dots,e_{n-1}\}$, of mean $\frac1nS_n$. The conditional expectation depends on $\psi$ only through $m$, so the tower property over independent passes gives $\E[(e_{pn},\bar\psi_{pn})]=M_n^p(e_0,e_0)$, and \eqref{eq:B4-rr-jensen} is $\E X^2\ge(\E X)^2$.
Closed form: by \eqref{eq:B4-rr-rec}, $S_{k+2}-(2-g)S_{k+1}+\varpi S_k=\frac{gk}nm$ with $S_0=0$, $S_1=e$; $m(k-n)$ is a particular solution, and $H_k:=S_k-m(k-n)$ solves the homogeneous recurrence with $H_0=nm$, $H_1=e+(n-1)m$, hence $H_k=H_1U_k-\varpi H_0U_{k-1}$ for $k\ge1$. Then $m'=\frac1nH_n$ and $\bar e_n=S_{n+1}-S_n=m+H_{n+1}-H_n$, which is \eqref{eq:B4-rr-closed}; expanding the determinant gives $\varpi(U_n^2-U_{n+1}U_{n-1})-\frac1nU_n$ and $U_n^2-U_{n+1}U_{n-1}=\varpi^{n-1}$ (induction on $n$ with the recurrence of $U$).
\end{proof}

\begin{proposition}[fluid limit of the pass map]\label{prop:B4-fluid}
Fix $c>0$ and let $\delta^2=\frac{c^2}4-c$, $h_1(u)=e^{-cu/2}\frac{\sinh\delta u}\delta$, $h_0(u)=e^{-cu/2}\big(\cosh\delta u+\frac c2\frac{\sinh\delta u}\delta\big)$ (for $c<4$, $\delta=iw$ with $w=\sqrt{c-c^2/4}$, $\frac{\sinh\delta u}\delta=\frac{\sin wu}w$ and $\cosh\delta u=\cos wu$; for $c=4$, $\frac{\sinh\delta u}{\delta}=u$).
\begin{enumerate}[label=(\roman*)]
\item $M_n(\frac cn)=M_\infty(c)+O(\frac1n)$ with (\texttt{rr\_fluid\_matrix})
\begin{equation}\label{eq:B4-fluid}
M_\infty(c)=\begin{pmatrix}h_1'(1)&1+h_0'(1)-h_1'(1)\\ h_1(1)&h_0(1)-h_1(1)\end{pmatrix},
\end{equation}
the map $(e,m)\mapsto(I'(1),I(1))$ of the equation $I''+cI'+cI=c\,m\,u$, $I(0)=0$, $I'(0)=e$, in which $I(u)$ is the running integral of the error over the pass.
\item $\det M_\infty(c)=e^{-c}-h_1(1)<0$ and $\tr M_\infty(c)=2e^{-c/2}\cosh\delta-h_1(1)$: the eigenvalues are real, of opposite signs, and the spectral radius $\nu(c)$ of $M_\infty(c)$ satisfies $\nu(c)<1$ for every $c>0$.
\item $\nu(c)=1-c+O(c^2)$ as $c\to0$, and $M_\infty(c)\to\begin{pmatrix}0&1-e^{-1}\\0&e^{-1}\end{pmatrix}$, $\nu(c)=e^{-1}\big(1+\frac{e-2}c\big)+O(c^{-2})$ as $c\to\infty$. Hence
\[
r^{\rm rr}(c):=-2\log\nu(c)=\lim_{n\to\infty}\big(-2\log\rho(M_n(\tfrac cn))\big),\qquad r^{\rm rr}(c)=2-\frac{2(e-2)}c+O(c^{-2})\ \xrightarrow[c\to\infty]{}\ 2
\]
($r^{\rm rr}$ is \texttt{rr\_exponent}).
\end{enumerate}
\end{proposition}
\begin{proof}
(i) With $u_k=\frac kn$ and $I_k=\frac1nS_k$, \eqref{eq:B4-rr-rec} is the explicit Euler scheme of step $\frac1n$ for the linear system $I'=\bar e$, $\bar e'=-c\,\bar e-c\,I+c\,m\,u$ on $[0,1]$, whose error at $u=1$ is $O(\frac1n)$ uniformly for $(e,m)$ in a bounded set (first-order convergence of Euler's method for a linear system with smooth coefficients on a compact interval). A particular solution is $m(u-1)$; $h_1$ and $h_0$ are the solutions of the homogeneous equation with $(h_1,h_1')(0)=(0,1)$ and $(h_0,h_0')(0)=(1,0)$ (the roots of $r^2+cr+c$ are $-\frac c2\pm\delta$), and $h_0'=-ch_1$; so $I=m(u-1)+m\,h_0+(e-m)h_1$, which gives \eqref{eq:B4-fluid}.
(ii) The Wronskian $h_0h_1'-h_0'h_1$ equals $e^{-cu}$, so $\det M_\infty=h_1'(h_0-h_1)-h_1(1+h_0'-h_1')=e^{-c}-h_1(1)$ at $u=1$, and $h_1'+h_0=2e^{-cu/2}\cosh\delta u$. If $c\ge4$, $h_1(1)=e^{-c/2}\frac{\sinh\delta}\delta\ge e^{-c/2}>e^{-c}$; if $c<4$, $\frac{\sin w}w\ge1-\frac{w^2}6\ge1-\frac c6$ and $e^{-c/2}\le1-\frac c2+\frac{c^2}8$ give $\frac{\sin w}w-e^{-c/2}\ge\frac c3(1-\frac{3c}8)>0$ for $c<\frac83$, while for $\frac83\le c<4$, $w\le1$ gives $\frac{\sin w}w\ge\sin1>e^{-4/3}\ge e^{-c/2}$. So $\det M_\infty<0$ and the characteristic polynomial $p(x)=x^2-\tr M_\infty\,x+\det M_\infty$ has two real roots of opposite signs; they lie in $(-1,1)$ iff $p(1)>0$ and $p(-1)>0$. Now $p(1)=1-2e^{-c/2}\cosh\delta+e^{-c}$ equals $(1-e^{-c/2+\delta})(1-e^{-c/2-\delta})>0$ for $c\ge4$ (as $\delta<\frac c2$) and is at least $(1-e^{-c/2})^2>0$ for $c<4$; and $p(-1)=1+e^{-c}+2e^{-c/2}\big(\cosh\delta-\frac{\sinh\delta}\delta\big)$, where the last bracket is nonnegative for $c\ge4$ ($\cosh\delta\ge\frac{\sinh\delta}\delta$ for real $\delta$) and at least $\cos1-1>-0.46$ for $c<4$ (it is then $\cos w-\frac{\sin w}w$ with $w\le1$); as $2e^{-c/2}\le2$, in all cases $p(-1)>1-0.92>0$.
(iii) Expand the closed forms. As $c\to0$: $\tr M_\infty=1-\frac{4c}3+O(c^2)$ and $\det M_\infty=-\frac c3+O(c^2)$, so the positive root is $1-c+O(c^2)$. As $c\to\infty$: $-\frac c2+\delta=-1-\frac1c+O(c^{-2})$ and $e^{-c/2-\delta}$ is exponentially small, so $e^{-c/2}\cosh\delta=\frac1{2e}(1-\frac1c)+O(c^{-2})$, $h_1(1)=\frac1{ec}+O(c^{-2})$, $h_1'(1)=O(c^{-1})$, $h_0(1)\to e^{-1}$, $h_0'(1)=-ch_1(1)\to-e^{-1}$; then $\tr M_\infty=e^{-1}(1-\frac2c)+O(c^{-2})$, $\det M_\infty=-\frac1{ec}+O(c^{-2})$ and the positive root is $\tr M_\infty-\det M_\infty/\tr M_\infty+O(c^{-2})=e^{-1}(1+\frac{e-2}c)+O(c^{-2})$. The limit in $n$ follows from (i) and the continuity of the eigenvalues.
\end{proof}

The first moments are not the whole story: the squared error is what an experiment measures. Its second moments close as well.

\begin{proposition}[exact second moments over a reshuffled pass, every $n\ge2$]\label{prop:B4-rr2}
In the setting of Proposition~\ref{prop:B4-rr}, let $v=\frac1n\sum_j\psi_j^2-m^2$ be the dispersion of the memory at the start of the pass and $\bar e_k=\E[e_k\mid e,\psi]$ the solution of \eqref{eq:B4-rr-rec}. There are constants $a_n,b_n,c_n,p_n$, depending on $(n,g)$ only, such that
\begin{equation}\label{eq:B4-rr2}
\begin{aligned}
\E[e'^2\mid e,\psi]&=\bar e_n^{\,2}+a_n\,v, & \E[m'^2\mid e,\psi]&=\big(\tfrac1nS_n\big)^2+c_n\,v,\\
\E[e'm'\mid e,\psi]&=\bar e_n\tfrac1nS_n+p_n\,v, & \E\Big[\tfrac1n\sum_j\psi_j'^{\,2}\,\Big|\,e,\psi\Big]&=\tfrac1n\sum_{k<n}\bar e_k^{\,2}+b_n\,v ,
\end{aligned}
\end{equation}
where $\psi'$ is the memory at the end of the pass and $S_n$ the sum of \eqref{eq:B4-rr-rec}. Since $\bar e_k$ and $S_n$ are linear in $(e,m)$, the moments $(A,B,C,P)$ of Proposition~\ref{prop:B4-iid}, taken at the ends of the passes, satisfy $(A,B,C,P)'=R_n(g)\,(A,B,C,P)$ for an explicit $4\times4$ matrix $R_n(g)$ (\texttt{rr\_moment\_matrix}; $e^2\mapsto A$, $em\mapsto P$, $m^2\mapsto C$, $v\mapsto B-C$). The constants are given by the covariance $\Sigma_k$ of $z_k=(e_k-\bar e_k,\ S_k-\E S_k,\ Y_k)$ for $v=1$, which obeys
\begin{equation}\label{eq:B4-rr2-cov}
\Sigma_{k+1}=F\Sigma_kF^\top+\kappa_kh^\top+h\kappa_k^\top+hh^\top,\qquad \kappa_k=-\frac{F\,\Sigma_k\,\mathbf e_3}{n-k},\qquad \Sigma_0=0,
\end{equation}
with $F=\begin{pmatrix}1-g&-\frac gn&\frac gn\\1&1&0\\0&0&1\end{pmatrix}$, $h=(g,0,1)^\top$, $\mathbf e_3=(0,0,1)^\top$: $a_n=(\Sigma_n)_{11}$, $c_n=\frac1{n^2}(\Sigma_n)_{22}$, $p_n=\frac1n(\Sigma_n)_{12}$, $b_n=\frac1n\sum_{k<n}(\Sigma_k)_{11}$.
\end{proposition}
\begin{proof}
Let $y_k=\psi_{\pi(k)}-m$ and $Y_k=y_1+\dots+y_k$: $(y_1,\dots,y_n)$ is a uniformly random ordering of a centred population of variance $v$, so $\E y_k^2=v$, $\E y_ky_l=-\frac v{n-1}$ ($k\ne l$), and, given $y_1,\dots,y_k$, the next term $y_{k+1}$ is uniform among the $n-k$ remaining values, whose sum is $-Y_k$: $\E[y_{k+1}\mid y_1,\dots,y_k]=-\frac{Y_k}{n-k}$. The sum of the unvisited slots after $k$ steps is $nm-km-Y_k$, so the step of the proof of Proposition~\ref{prop:B4-rr} reads
\[
e_{k+1}=(1-g)e_k-\frac gnS_k+\frac{gk}nm+g\,y_{k+1}+\frac gnY_k .
\]
Subtracting \eqref{eq:B4-rr-rec}, $z_{k+1}=Fz_k+h\,y_{k+1}$, where $z_k$ is a linear function of $y_1,\dots,y_k$ with coefficients depending on $(n,g)$ only; hence $\E z_k=0$, $\E[z_ky_{k+1}]=-\frac1{n-k}\E[z_kY_k]$, and the covariance of $z_k$ is $v\,\Sigma_k$ with \eqref{eq:B4-rr2-cov}. Then $e'=e_n=\bar e_n+(z_n)_1$, $m'=\frac1n(S_n+(z_n)_2)$ and $\frac1n\sum_j\psi_j'^2=\frac1n\sum_{k<n}e_k^2$ give \eqref{eq:B4-rr2}, the cross terms vanishing because $\E z_k=0$. The conditional moments depend on $\psi$ through $(m,v)$ only, and a fresh permutation is independent of the past, so the tower property gives the linear map on $(A,B,C,P)$.
\end{proof}

\begin{proposition}[mean-square rate of the reshuffled order: the two regimes]\label{prop:B4-rr-ms}
Let $r^{\rm rr}_n(c):=-\log\rho(R_n(\frac cn))$ be the exact exponent per pass of $\E e^2$ at finite $n$ (\texttt{rr\_ms\_exponent}; by \eqref{eq:B4-rr-jensen} it is at most $-2\log\rho(M_n(\frac cn))$).
\begin{enumerate}[label=(\roman*)]
\item \emph{Fixed $c$ (proved).} $a_n,b_n,c_n,p_n=O(\frac1n)$ as $n\to\infty$ at fixed $c$, and
\[
\lim_{n\to\infty}r^{\rm rr}_n(c)=-2\log\nu(c)=r^{\rm rr}(c):
\]
the expected squared error of reshuffled SAGA decays at the rate of its squared mean.
\item \emph{Fixed $g\in(0,2)$ (proved).} As $n\to\infty$ at fixed $g$, $a_n$ and $b_n$ tend to the fluctuation factor $\varphi(g)=\frac g{2-g}$, $c_n$ and $p_n$ to $0$, $M_n(g)$ to $\begin{pmatrix}0&1-e^{-1}\\0&e^{-1}\end{pmatrix}$, the spectrum of $R_n(g)$ to $\{\varphi(g),e^{-2},0,0\}$, and
\[
\lim_{n\to\infty}r^{\rm rr}_n(ng)=\min\Big\{2,\ \log\frac{2-g}g\Big\}=:r^{\rm rr}_{\rm fix}(g)
\]
(\texttt{rr\_exponent\_fixed\_g}). The second branch is the smaller one for $g>\frac2{1+e^2}=0.2384$; it equals $\log7=1.946$ at $g=\frac14$ and is negative for $g>1$: there the mean square \emph{diverges}, although the mean converges at two nats per pass.
\end{enumerate}
\end{proposition}
\begin{proof}
(i) Write $e_k-\bar e_k=\sum_lW_{k,l}\,y_l$. From the recursion above, $W_{k+1,l}=(1-g)W_{k,l}-\frac gn\sum_{j<k}W_{j,l}+\frac gn\ind[l\le k]+g\,\ind[l=k+1]$, so $u_k:=\max_{j\le k,\,l}\abs{W_{j,l}}$ satisfies $u_{k+1}\le\max\{u_k(1+\frac{ck}{n^2})+\frac c{n^2},\ \frac cn\}$ and therefore $u_n\le\frac{2c}ne^{c/2}=:\frac{C}n$. For weights $w$ with $\abs{w_l}\le\frac Cn$, $\E(\sum_lw_ly_l)^2=v\big(\frac n{n-1}\sum_lw_l^2-\frac1{n-1}(\sum_lw_l)^2\big)\le\frac{2C^2}nv$; this applies to $e_k-\bar e_k$ and to $\frac1n(S_n-\E S_n)$ (an average of such terms), and the Cauchy--Schwarz inequality bounds the mixed term: $a_n,b_n,c_n,\abs{p_n}\le\frac{2C^2}n$. In the coordinates $(A,P,C,v)$ with $v=B-C$, $R_n$ therefore has the block form $\begin{pmatrix}{\rm Sym}^2M_n&O(\frac1n)\\ d_n&O(\frac1n)\end{pmatrix}$, where ${\rm Sym}^2M_n$ is the action of $M_n(\frac cn)$ on the quadratic monomials $(e^2,em,m^2)$ and the row $d_n$ (the dispersion of $\bar e_0,\dots,\bar e_{n-1}$ as a quadratic form in $(e,m)$) is bounded. As $M_n(\frac cn)\to M_\infty(c)$ (Proposition~\ref{prop:B4-fluid}(i)), $R_n$ converges to a block-triangular matrix with spectrum $\{\nu_1^2,\nu_1\nu_2,\nu_2^2,0\}$, $\nu_{1,2}$ the eigenvalues of $M_\infty(c)$; eigenvalues are continuous, so $\rho(R_n(\frac cn))\to\nu(c)^2$.
(ii) Take $n>\frac4g$, so that the roots $r_\pm=1-\frac g2\pm\sqrt{\frac{g^2}4-\frac gn}$ of $x^2-(2-g)x+\varpi$ are real: $r_+=1-\frac1n+O(n^{-2})$, $r_-=1-g+\frac1n+O(n^{-2})$, $r_+r_-=\varpi$, $r_+-r_-\to g$, and $\abs{r_-}\le\varrho$ for some $\varrho<1$ and every large $n$ (this is where $0<g<2$ is used). Below, $C$ is a constant depending on $g$ only.
\emph{Means.} $U_k=\frac{r_+^k-r_-^k}{r_+-r_-}$ is bounded, $r_+^n\to e^{-1}$, $r_-^n\to0$; hence in \eqref{eq:B4-rr-closed} $U_{n+1}-U_n\to0$, $\frac1nU_n\to0$, $n(U_{n+1}-U_n)\to-\frac1{ge}$, $n\varpi(U_n-U_{n-1})\to-\frac{1-g}{ge}$ and $U_n-\varpi U_{n-1}=\frac{r_+^n(1-r_-)-r_-^n(1-r_+)}{r_+-r_-}\to e^{-1}$, which gives the limit of $M_n(g)$. Moreover, by the proof of Proposition~\ref{prop:B4-rr}, $\bar e_k=m+H_{k+1}-H_k=e\,u_k+m\big(1-u_k+n(r_+-1)(1-r_-)U_k\big)$ with $u_k=U_{k+1}-U_k$ (use $(U_{k+1}-U_k)-\varpi(U_k-U_{k-1})=(r_+-1)(1-r_-)U_k$), and $\abs{u_k}\le C(\varrho^k+\frac1n)$: the coefficients of $\bar e_k$ are bounded uniformly in $k\le n$, and that of $e$ is summable up to $O(1)$. So the row $d_n$ of (i), the quadratic form $\frac1n\sum_{k<n}\bar e_k^{\,2}-(\frac1nS_n)^2$ in $(e,m)$, is bounded, and its coefficient of $e^2$, $\frac1n\sum_{k<n}u_k^2-(\frac1nU_n)^2$, is $O(\frac1n)$.
\emph{Fluctuations.} The recursion of the weights in (i) does not depend on $l$: $W_{k,l}=w_{k-l}$ for $k\ge l$ (and $0$ for $k<l$), with $w_0=g$, $w_{j+1}=(1-g)w_j-\frac gn(t_j-1)$, $t_j=\sum_{i<j}w_i$. The sequence $\tau_j=t_j-1$ solves $\tau_{j+2}=(2-g)\tau_{j+1}-\varpi\tau_j$ with $\tau_0=-1$, $\tau_1=g-1$, so $\tau_j=\alpha r_+^j+\beta r_-^j$ with $\alpha=\frac{g-1+r_-}{r_+-r_-}=O(\frac1n)$ and $\beta=-1-\alpha$, and
\[
w_j=\tau_{j+1}-\tau_j=g_n\,r_-^{\,j}+\alpha(r_+-1)r_+^{\,j},\qquad g_n=\beta(r_--1)=g+O(\tfrac1n),\qquad \abs{\alpha(r_+-1)r_+^{\,j}}\le\frac C{n^2},\qquad \abs{t_j}\le C .
\]
By the variance formula of (i), exact for every weight vector, $(\Sigma_k)_{11}=\frac n{n-1}\sum_{j<k}w_j^2-\frac1{n-1}t_k^2=\varphi(g)\big(1-(1-g)^{2k}\big)+O(\frac1n)$ uniformly in $k\le n$, because $\sum_{j\ge0}g^2(1-g)^{2j}=\frac{g^2}{1-(1-g)^2}=\varphi(g)$: $a_n=(\Sigma_n)_{11}\to\varphi$ and $b_n=\frac1n\sum_{k<n}(\Sigma_k)_{11}\to\varphi$. Next $S_n-\E S_n=\sum_{l<n}t_{n-l}\,y_l$ has bounded weights, so $(\Sigma_n)_{22}\le Cn$, $c_n=O(\frac1n)$, and $\abs{p_n}\le\sqrt{a_nc_n}=O(n^{-1/2})$.
\emph{Spectrum.} In the coordinates $(A,P,C,v)$ of (i), $R_n(g)$ has the blocks ${\rm Sym}^2M_n(g)$, the column $(a_n,p_n,c_n)^\top$, the row $d_n$ and the entry $b_n-c_n$. It is bounded, and each of its limit points is
\[
\begin{pmatrix}0&0&(1-e^{-1})^2&\varphi\\0&0&(1-e^{-1})e^{-1}&0\\0&0&e^{-2}&0\\0&d_P&d_C&\varphi\end{pmatrix},
\]
whose characteristic polynomial is $x^2(x-e^{-2})(x-\varphi)$ whatever $d_P$ and $d_C$ (expand along the third row, then along the second). Eigenvalues are continuous, so $\rho(R_n(g))\to\max\{e^{-2},\varphi(g)\}$. (In fact $d_n\to(0,0,\frac2e-\frac12-\frac3{2e^2})$, the dispersion of $1-e^{-u}$ over $u\in[0,1]$: the limit of the mean path.)
\emph{Check on the exact matrix} (\texttt{rr\_ms\_exponent}): $r^{\rm rr}_n(ng)=1.9994,\ 1.9448,\ 1.0984,\ -0.5108$ at $n=10^4$ for $g=0.2,\ \frac14,\ \frac12,\ 1.25$, against $2,\ 1.9459,\ 1.0986,\ -0.5108$; the differences decrease like $\frac1n$ between $n=10^3$ and $10^5$.
% B4b_summary.json: note.R3_fixed_g, fixed_g.by_g.<g>.{rr_exact, rr_limit, rr_mean}
\end{proof}

\begin{remark}[status of (R3), numerical values, and the best step]\label{rem:B4-rr-status}
\emph{Proved:} the mean error of reshuffled SAGA after $p$ passes is given exactly by $M_n(g)^p$ and its second moments by $R_n(g)^p$, at every $n\ge2$; at fixed $c$ the expected squared error decays at $r^{\rm rr}(c)$ in the limit $n\to\infty$, and at a fixed step $g\in(0,2)$ at $\min\{2,\log\frac{2-g}g\}$; the limiting first-moment map is stable for every $c$. The limiting picture is transparent for large $c$: the error follows the memory, $I'+I=mu$, and the mean of the memory is multiplied by $e^{-1}$ at each pass, every slot being replaced once.
\emph{Not proved:} that a \emph{typical} run decays at the rate of $\E e^2$. The latter is supported by the simulations of the Outcome, and is expected because the dispersion $v$ is an average over $n$ slots; for i.i.d.\ sampling it is false (Section~\ref{sec:B4-typical}).
\emph{Values} (\texttt{rr\_exponent}): $r^{\rm rr}(c)=0.445$, $0.777$, $1.187$, $1.537$, $1.752$, $1.863$, $1.923$, $1.959$ at $c=\frac14,\frac12,1,2,4,8,16,32$; on a grid of $4001$ values of $c$ in $[10^{-3},10^5]$ it is increasing and smaller than $\min\{2c,2\}$, and it equals $1$ at $c=0.730$ (numerical); $c\,(2-r^{\rm rr}(c))=1.4326$ and $1.4362$ at $c=10^3$ and $10^4$, against $2(e-2)=1.4366$. \emph{Exact finite-$n$ values:} for $n=200$, $c=4$, after six passes from the uniform memory the mean is $5.4165\cdot10^{-3}$ and the root mean square $5.4527\cdot10^{-3}$, $0.67\%$ above it; $r^{\rm rr}_n(c)$ and the exponent of the squared mean agree to within $2\cdot10^{-6}$ at $n=1000$ for $c=1,4,16,32$ ($1.18802$, $1.75337$, $1.92400$, $1.95961$), but for $n=100$, $c=32$ ($g=0.32$) the mean square decays at $1.6030$ whereas the squared mean decays at $1.9685$: the dispersion branch of Proposition~\ref{prop:B4-rr-ms}(ii), $\log\frac{2-g}g=1.658$ in the limit.
\emph{Best step.} At finite $n$ the exponent $r^{\rm rr}_n(c)$ increases with $c$ along the first branch and decreases along the second: it is largest where they cross, near $g=0.24$, with the values $1.952$ at $c=22.3$ for $n=100$, $1.995$ at $c=237$ for $n=1000$ and $1.9995$ at $c=2.4\cdot10^3$ for $n=10^4$ (\texttt{rr\_ms\_exponent}, bounded search). Two nats per pass are approached only when $n$ is large and the step is chosen at $\gamma\lambda\approx0.24$ on the stiffest mode, which the textbook value $\frac14$ nearly is.
% B4b_summary.json: note.R3_exponent, note.R3_crossing_r_equals_1, note.R3_grid_4001, note.R3_large_c,
% note.R3_n200_c4_after_6_passes, note.R3_mean_square_vs_mean, fixed_g.rr_best
\end{remark}

\subsection{Comparison of the three orders}\label{sec:B4-comparison}

\begin{center}
\footnotesize
\begin{tabular}{@{}lcccccccc@{}}
\toprule
$c=\gamma n\lambda$ & $\frac14$ & $\frac12$ & $1$ & $2$ & $4$ & $8$ & $16$ & $32$ \\
\midrule
i.i.d., expectation: $r^{\rm iid}=\min\{2c,1\}$ & $0.5$ & $1$ & $1$ & $1$ & $1$ & $1$ & $1$ & $1$ \\
fixed order: $r^{\rm cyc}$ & $0.437$ & $0.756$ & $1.149$ & $1.497$ & $0.702$ & $0.194$ & $0.011$ & $-0.025$ \\
reshuffled: $r^{\rm rr}$ & $0.445$ & $0.777$ & $1.187$ & $1.537$ & $1.752$ & $1.863$ & $1.923$ & $1.959$ \\
\bottomrule
\end{tabular}
\end{center}
\noindent Limits $n\to\infty$ at fixed $c$ of the squared-error exponents per pass of one mode (Propositions~\ref{prop:B4-iid-limit}, \ref{prop:B4-cyclic-limit}, \ref{prop:B4-fluid}, \ref{prop:B4-rr-ms}(i); for the first and third rows the exponent of the expected squared error). At a fixed step $g=\gamma\lambda$ instead (Propositions~\ref{prop:B4-iid-fixed}, \ref{prop:B4-rr-ms}(ii); a fixed order is unstable for $n$ large at every fixed $g>0$, since $c=ng>c_\ast$):

\begin{center}
\footnotesize
\begin{tabular}{@{}lcccccccc@{}}
\toprule
$g=\gamma\lambda$ & $0.05$ & $0.1$ & $0.2$ & $\frac14$ & $0.32$ & $\frac12$ & $0.9$ & $1.1$ \\
\midrule
i.i.d., expectation: $\frac{2(1-g)}{2-g}$ & $0.974$ & $0.947$ & $0.889$ & $0.857$ & $0.810$ & $0.667$ & $0.182$ & $-0.222$ \\
reshuffled: $\min\{2,\log\frac{2-g}g\}$ & $2$ & $2$ & $2$ & $1.946$ & $1.658$ & $1.099$ & $0.201$ & $-0.201$ \\
\bottomrule
\end{tabular}
\end{center}
\noindent Four conclusions, with their status.
\begin{enumerate}[label=(\arabic*),nosep]
\item \emph{The ceiling belongs to sampling with replacement.} For i.i.d.\ indices the expected squared error decays at $\min\{2c,1\}$ at fixed $c$ and at $\frac{2(1-g)}{2-g}<1$ at a fixed step: at most one nat per pass (proved). For reshuffled indices it decays at $r^{\rm rr}(c)$, which exceeds $1$ for $c>0.730$ (numerical) and tends to $2$ (proved): the memory rate is doubled, not removed, since a slot refreshed once per pass still holds an error one pass old.
\item \emph{Reshuffling is not uniformly better.} $r^{\rm rr}(c)<2c$ (numerical): for $c<0.730$ the reshuffled exponent is below the i.i.d.\ one ($0.445$ against $0.5$ at $c=\frac14$); under reshuffling the direction is no longer an unbiased estimate of the gradient, the slot sampled being uniform among those not yet visited in the pass. With several modes the squared error is governed by the slowest of them, which may be a soft mode ($c_{\min}=\gamma n\lambda_1$, $r^{\rm rr}$ being increasing) or the stiffest one through its dispersion ($g_{\max}=\gamma\lambda_d$): the exponent is $\min\{r^{\rm rr}(c_{\min}),\log\frac{2-g_{\max}}{g_{\max}}\}$ up to $O(\frac1n)$. Reshuffling pays when $\gamma n\mu\gtrsim1$, i.e.\ for the textbook step when $\kappa\lesssim n/4$, the regime in which the i.i.d.\ rate is the memory rate.
\item \emph{The step matters differently for the two orders.} At the textbook step the stiffest mode has $g=\frac14$: in expectation the i.i.d.\ order then gives $\frac67$ nat per pass and the reshuffled order $\log7=1.946$, a ratio $2.27$ (the first proved, the second verified on the exact matrix). Larger steps help neither order beyond $g\approx0.24$, and both diverge in mean square beyond $g=1$; for the reshuffled order this instability is invisible on the mean, which contracts for every step.
\item \emph{A fixed order is not a substitute.} It is within $4\%$ of reshuffling at $c\le2$ in the first table and unstable beyond $c_\ast$ (proved in the limit), with a stability condition on the \emph{stiffest} mode, $\gamma n\lambda_d<c_\ast(n)$: the step must be $O(\frac1{nR^2})$, then $c_{\min}<c_\ast/\kappa$ and, since $r^{\rm cyc}(c)<2c$ at every value computed, the exponent is at most about $2c_\ast/\kappa$ per pass. The textbook step is excluded for the values of $n\ge52$ that were checked (Remark~\ref{rem:B4-visible}(b)). This is the quantitative form, in the model, of an observation as old as these methods (Section~\ref{sec:B4-related}).
\end{enumerate}
These are statements on expectations (first and third conclusions) or on a deterministic sequence (fourth). What a single i.i.d.\ run does is another matter (Section~\ref{sec:B4-typical}).

\subsection{The lower bounds of the note on lower bounds under reshuffling}\label{sec:B4-lower}

On the instance \eqref{eq:instance} with $\zeta=\zeta_Q$ ((S) in the report; $Q=1+\frac{\kappa-1}n$, $q=\frac{\sqrt Q-1}{\sqrt Q+1}$, $\theta_0=0$), every bound of Section~\ref{sec:lb:instance} is the tail bound \eqref{ineq:tail} combined with a statement on the numbers $\tau_j$ of queries to each component. The sampling enters at one place only: in Theorem~\ref{thm:instance} and Corollaries~\ref{cor:count} and~\ref{cor:sweep}, ``under (S1), $\tau_j(t)\sim\mathrm{Binomial}(t,p_j)$''. This is the hypothesis that a reshuffled, shuffled-once or cyclic method leaves; such a method is still in class (S2) of Definition~\ref{def:span}, so Lemma~\ref{lem:support} and Proposition~\ref{prop:anyseq} apply to it.

\begin{proposition}[floors on the instance, by order]\label{prop:B4-floors}
Run SAGA with $z_i^{(0)}=\fip(\theta_0)$, $\theta_0=0$, on the instance \eqref{eq:instance} with $\zeta=\zeta_Q$, for $p$ passes ($G=n(p+1)$ gradient evaluations), $p+1\le m$.
\begin{enumerate}[label=(\alph*)]
\item I.i.d.\ indices: $\dfrac{\E\norm{\theta_{pn}-\etastar}^2}{\norm{\etastar}^2}\ge\dfrac{q^2(1-\frac{1-q^2}n)^{pn}-q^{2m}}{1-q^{2m}}$ (this is \eqref{ineq:expected-sweep}); the floor loses $-n\log(1-\frac{1-q^2}n)\le\frac{1-q^2}{1-1/n}$ nats per pass.
\item Any order that visits every index once per pass: pathwise, $\dfrac{\norm{\theta_{pn}-\etastar}^2}{\norm{\etastar}^2}\ge\dfrac{q^{2(p+1)}-q^{2m}}{1-q^{2m}}$, which is the floor \eqref{ineq:jensen} of Proposition~\ref{prop:anyseq} at $G=n(p+1)$, with equality in Jensen's inequality; it loses $2\log\frac1q$ nats per pass.
\end{enumerate}
For $n\ge2$ and $w=1-q^2\in(0,1)$, $-n\log(1-\frac wn)<-\log(1-w)=2\log\frac1q$ (compare the power series term by term), so the floor (b) is below the floor (a) for every $p\ge1$: a method without replacement that goes below the i.i.d.\ bound of Corollary~\ref{cor:sweep} contradicts nothing, and its gradient count remains bounded below by \eqref{eq:anyseqcount}.
\end{proposition}
\begin{proof}
(a) is Corollary~\ref{cor:sweep} at $t=pn$, and $-\log(1-x)\le\frac x{1-x}$. (b) By Lemma~\ref{lem:support}(ii), $N_j(\theta_{pn})\le1+\tau_j(pn)=1+p$ for every $j$; apply Proposition~\ref{prop:tail} with $\tau_j=p+1\le m$.
\end{proof}

\paragraph{What the model predicts for the ratio of gradient counts.}
The floor (a) decays at essentially one nat per pass when $\kappa\ll n$ (for $n=256$: $q=0.0229$ and $1.0014$ nats per pass at $\kappa/n=0.1$; $q=0.0711$ and $0.9969$ at $\kappa/n=\frac13$; against $2\log\frac1q=7.55$ and $5.29$ for the floor (b)): it is the lower-bound counterpart of the memory ceiling, a block never sampled playing the role of a slot never refreshed. If the squared error of a reshuffled method decays at $r$ nats per pass, its count is $n+\frac nr\log\frac C\varepsilon$, against $n+(\frac n{1-q^2}-1)\log\frac{q^2}{2\varepsilon}$ for the bound of Corollary~\ref{cor:sweep}: the ratio tends to $\frac{1-q^2}r$ as $\varepsilon\to0$ (up to a factor $1+O(\frac1n)$), and differs from it at moderate accuracy because of the constants inside the logarithms. In the solvable model $r<2$ for reshuffled SAGA (at every value computed), with $r\to2$ at large $\gamma n\lambda_{\min}$. \emph{If} this ceiling carried over to the instance, the ratio would be at least $\frac{1-q^2}2$, i.e.\ reshuffling would at best halve the i.i.d.\ bound asymptotically. This is an analogy, not a statement on the instance, which does \emph{not} satisfy Assumption~\ref{ass:B4-model} (the Hessians of the components live on different blocks): what is proved there for an order without replacement is only the floor (b), which leaves room for the factor $2\log\frac1q/(1-q^2)$ ($7.55$ and $5.32$ for the two values of $q$ above) that full-gradient schemes such as the SVRG of Theorem~\ref{thm:HLOY}(b) exploit. The rate $r$ on the instance must be measured. A second caution: the bound (a) concerns an expectation that is carried by rare runs (Section~\ref{sec:B4-typical}), whereas (b) is pathwise.
% B4b_summary.json: note.floors_n256

\subsection{Expectation against typical runs under i.i.d.\ sampling}\label{sec:B4-typical}

Under i.i.d.\ sampling the expectation and the typical run must be kept apart, in the model as on the hard instance. Write $x_t=(e_t,\psi^{(t)})\in\R^{n+1}$ for the state of the scalar recursion \eqref{eq:B4-scalar}: $x_t=A_{i(t)}x_{t-1}$ for $n$ fixed matrices $A_1,\dots,A_n$.
\begin{enumerate}[label=(\alph*),nosep]
\item \emph{(Proved.)} Let $N_t$ be the number of slots (or blocks) never sampled during the first $t$ iterations. Then $\E N_t=n(1-\frac1n)^t\le ne^{-t/n}$ and $\Prob[N_t\ge1]\le ne^{-t/n}$: after $\log n+a$ passes a run has no such slot with probability at least $1-e^{-a}$. In an average over $s$ independent runs, the expected number of (run, slot) pairs never sampled after $p$ passes is at most $s\,n\,e^{-p}$, below $1$ as soon as $p>\log(sn)$ ($9.0$ passes for $s=32$ runs and $n=256$).
\item \emph{(Proved: a typical exponent exists.)} The matrices $A_{i(t)}$ are i.i.d.\ and bounded, so by the theorem of Furstenberg and Kesten \cite{furstenberg1960products} $\frac1t\log\norm{A_{i(t)}\cdots A_{i(1)}}$ converges almost surely to a constant $\Lambda_n(g)$, the top Lyapunov exponent. We call $r^{\rm typ}_n(c):=-2n\Lambda_n(\frac cn)$ the \emph{typical exponent}: it is what the slope of $\log\norm{x_t}^2$ of a single long run estimates.
\item \emph{(Proved: a run is at least as fast as the expectation.)} $\E\norm{x_t}^2=A_t+nB_t\le K\rho_n^t$ with $\rho_n$ the spectral radius of $T_n(g)$. By Markov's inequality $\Prob[\norm{x_t}^2>\rho_n^te^{\epsilon t/n}]\le Ke^{-\epsilon t/n}$, which is summable in $t$: by the Borel--Cantelli lemma, almost surely $\limsup_t\frac nt\log\norm{x_t}^2\le-r^{\rm iid}_n(c)+\epsilon$ for every $\epsilon>0$. Hence $r^{\rm typ}_n(c)\ge r^{\rm iid}_n(c)$ whenever the run started from the uniform memory realizes the top exponent, and in every case almost every run decays at least as fast as the expectation.
\item \emph{(Heuristic: why it is strictly faster, and by how much.)} In (R1) with $c>\frac12$ the expectation \eqref{eq:B4-plateau} is of the order of $\frac{e_0^2}{n^2}$ times the expected number $ne^{-t/n}$ of never-refreshed slots: it is carried by runs in which a slot of the initial memory survived until shortly before $t$, which are rare for $t/n>\log n$. For a run that has forgotten $e_0$ the same mechanism acts on the oldest slots it still holds. Suppose $\norm{x_t}^2$ decays at $r$ nats per pass. A slot of age $a$ passes holds an error $e^{ra/2}$ times the current one, and ages are exponential: the stored values have a power-law tail of index $\frac2r\in(1,2)$ for $1<r<2$. When $c\gg1$ the error is slaved to the memory and equals the centred mean of the $N=\frac nc$ slots sampled during its response time; for a tail of index $\frac2r$ this mean is of relative size $N^{r/2-1}$ times $\frac2{2-r}$ (the largest value sampled, the centring being dominated by larger values that are typically not sampled). Self-consistency requires this to be $1$:
\begin{equation}\label{eq:B4-typical}
r^{\rm typ}_n(c)\approx2-\delta,\qquad \delta\,\log\frac nc=2\log\frac2\delta\qquad(1\ll c\ll n),
\end{equation}
and $r^{\rm typ}\approx2c$ when $2c$ is smaller (\texttt{iid\_typical\_heuristic}). This is not a proof.
\end{enumerate}
\emph{Measured values.} The exponent of single runs (median over $12$ seeds, $260$ passes, fitted from pass $60$ on; experiment B4b) is, against \eqref{eq:B4-typical} in parentheses:
% B4b_summary.json: typical.n=<n>.<c>.{median, heuristic, expectation, rr}, typical.heuristic_error
\begin{center}
\footnotesize
\begin{tabular}{@{}lcccc@{}}
\toprule
$c$ & $n=10^2$ & $n=10^3$ & $n=10^4$ & $n=10^5$ \\
\midrule
$1$ & $1.55$ ($1.44$) & $1.61$ ($1.56$) & $1.67$ ($1.63$) & $1.73$ ($1.68$) \\
$2$ & $1.49$ ($1.39$) & $1.63$ ($1.53$) & $1.67$ ($1.61$) & $1.77$ ($1.67$) \\
$4$ & $1.38$ ($1.33$) & $1.56$ ($1.50$) & $1.65$ ($1.59$) & $1.71$ ($1.65$) \\
$8$ & $1.27$ ($1.24$) & $1.48$ ($1.46$) & $1.61$ ($1.57$) & $1.67$ ($1.64$) \\
$16$ & $1.12$ ($1.11$) & $1.40$ ($1.41$) & $1.53$ ($1.54$) & $1.63$ ($1.62$) \\
$32$ & $0.92$ ($0.92$) & $1.35$ ($1.35$) & $1.50$ ($1.51$) & $1.60$ ($1.60$) \\
\midrule
reshuffled, exact, $c=32$ & $1.60$ & $1.96$ & $1.96$ & $1.96$ \\
i.i.d., expectation, exact, $c=32$ & $0.81$ & $0.98$ & $1.00$ & $1.00$ \\
\bottomrule
\end{tabular}
\end{center}
\noindent A typical run is faster than the expectation in every cell of the table, by $0.11$ ($n=100$, $c=32$) to $0.77$ nats per pass, and its exponent increases with $n$ at every $c$. The heuristic \eqref{eq:B4-typical} is within $0.043$ of the measured median (median deviation $0.013$) on the $9$ cells with $c\ge8$ and $\frac cn\le0.05$, and less accurate at small $c$, where the error is not slaved (deviations up to $0.11$, median $0.06$, for $c\le4$, always below the measured value). The spread over the seeds (largest minus smallest) is $0.04$ to $0.20$. At a fixed step the typical exponent \emph{decreases} with $g$: for $n=10^3$ it is $1.24$, $1.13$ and $0.96$ at $c=64,125,250$, the last one being the textbook value $g=\frac14$, where the expectation decays at $0.857$ and the reshuffled order at $1.935$: there reshuffling is $2.02$ times faster than a typical i.i.d.\ run. At equal $c\le32$ the ratio of the reshuffled exponent to the typical i.i.d.\ one is below $1$ at $c=1$ (and at $c=2$ for $n\ge10^3$) and reaches $1.74$, $1.45$, $1.31$ and $1.22$ at $c=32$ for $n=10^2,10^3,10^4,10^5$: the advantage of reshuffling over typical i.i.d.\ runs shrinks slowly as $n$ grows, as the conjecture below suggests.

\emph{Conjecture.} At fixed $c\ge1$, $r^{\rm typ}_n(c)\to2$ as $n\to\infty$, with $2-r^{\rm typ}_n(c)$ of order $\frac{\log\log n}{\log n}$ (the solution of \eqref{eq:B4-typical} is $\delta\sim\frac{2\log\log n}{\log n}$); at fixed $c<1$, $r^{\rm typ}_n(c)\to2c$. If so, the two orders have the \emph{same} limiting rate for typical runs at large $c$ and differ by the speed at which it is approached: $\frac{2(e-2)}c$ below $2$ for reshuffling, whatever $n$, against a logarithmically small quantity in $n$ for i.i.d.\ sampling, at a much smaller best step ($c$ of order one against $g\approx0.24$).

Hence, beyond about $\log n$ passes, $\E e_t^2$ is carried by rare runs, a typical run decays faster than the expectation, and a seed average over few runs estimates the typical behaviour, not the expectation: it may fall below a lower bound proved for the expectation, such as Proposition~\ref{prop:B4-floors}(a), without contradicting it. The comparison ``two nats against one'' of Section~\ref{sec:B4-comparison} is a comparison of the reshuffled order with the \emph{expectation} under i.i.d.\ sampling; against typical i.i.d.\ runs the gain of reshuffling is the smaller one of the table above. For the reshuffled order the distinction does not arise at small $g$: its mean, its mean square and its simulated runs decay at the same rate (Remark~\ref{rem:B4-rr-status}, Outcome).

\subsection{Beyond identical Hessians: the effective curvature, and the orthogonal extreme}\label{sec:B4-beyond}

Everything above rests on Assumption~\ref{ass:B4-model}. Two computations show what survives without it. Let now $f_i(\theta)=\frac12\theta^\top H_i\theta-b_i^\top\theta$ with symmetric positive semi-definite $H_i$ and $\bar H=\frac1n\sum_iH_i$ positive definite (or let $H_i=f_i''(\etastar)$ for smooth $f_i$, as a local description). The computation of Lemma~\ref{lem:B4-reduction} gives
\begin{equation}\label{eq:B4-general}
e_t=e_{t-1}-\gamma\Big(H_{i(t)}\big(e_{t-1}-\psi^{(t-1)}_{i(t)}\big)+\frac1n\sum_jH_j\psi_j^{(t-1)}\Big),\qquad\psi^{(t)}_{i(t)}=e_{t-1}.
\end{equation}

\paragraph{The effective curvature.} The fixed-step laws of Propositions~\ref{prop:B4-iid-fixed} and~\ref{prop:B4-rr-ms}(ii) come from one mechanism: an error slaved to the memory turns a dispersion of the stored points into a dispersion of the refreshed ones, smaller by the fluctuation factor. For \eqref{eq:B4-general}, let the stored points be $\psi_j=\bar\psi+\xi_j$ with centred, uncorrelated $\xi_j$ of covariance $\Xi$. Around its mean the error follows $\delta e^+=(I-\gamma H_i)\delta e+\gamma H_i\xi_i$ (up to the terms that vanish in the slaved regime), whose stationary covariance $\Sigma$ solves
\begin{equation}\label{eq:B4-fluct}
\Sigma=\frac1n\sum_i(I-\gamma H_i)\Sigma(I-\gamma H_i)+\gamma^2\mathcal K(\Xi),\qquad \mathcal K(\Xi):=\frac1n\sum_iH_i\Xi H_i ,
\end{equation}
and the refreshed slots inherit $\Sigma$.

\begin{proposition}[effective curvature]\label{prop:B4-Leff}
Let $\mathcal L(\Xi):=\bar H\Xi+\Xi\bar H$ on symmetric matrices and
\[
\tilde L:=2\,\lmax\big(\mathcal L^{-1}\mathcal K\big)
\]
(\texttt{effective\_curvature}). The spectrum of $\mathcal L^{-1}\mathcal K$ is real and non-negative, $\lmin(\bar H)\le\tilde L\le\max_i\lmax(H_i)$ (i.e.\ $\mu\le\tilde L\le R^2$), and $\tilde L=\lmax(H)$ when all $H_i$ equal $H$. For $0<\gamma\tilde L<2$ the linear map $\Xi\mapsto\Sigma$ defined by \eqref{eq:B4-fluct} has spectral radius
\[
\varphi(\gamma\tilde L)=\frac{\gamma\tilde L}{2-\gamma\tilde L},
\]
the fluctuation factor at $g=\gamma\tilde L$.
\end{proposition}
\begin{proof}
For the trace inner product $\mathcal L$ is self-adjoint positive definite ($\tr(X\mathcal L(X))=2\tr(X\bar HX)>0$) and $\mathcal K$ is self-adjoint positive semi-definite ($\tr(X\mathcal K(X))=\frac1n\sum_i\norm{H_i^{1/2}XH_i^{1/2}}_F^2$), so $\mathcal L^{-1}\mathcal K$ is self-adjoint for $\langle X,Y\rangle_{\mathcal L}=\tr(X\mathcal L(Y))$ with non-negative eigenvalues. Both $\mathcal K$ and $\mathcal L^{-1}$ (the solution of a Lyapunov equation with positive semi-definite right-hand side) preserve the cone of positive semi-definite matrices, so the largest eigenvalue $\frac{\tilde L}2$ has an eigenvector $\Xi\succeq0$ (Krein--Rutman): $\mathcal K(\Xi)=\frac{\tilde L}2\mathcal L(\Xi)$. Taking traces, $\frac1n\sum_i\tr(H_i^2\Xi)=\tilde L\,\tr(\bar H\Xi)$; the left-hand side is at most $\max_i\lmax(H_i)\tr(\bar H\Xi)$ and, since $\frac1n\sum_iH_i^2\succeq\bar H^2$, at least $\lmin(\bar H)\tr(\bar H\Xi)$, with $\tr(\bar H\Xi)>0$. For $H_i\equiv H$ with eigenpairs $(\lambda_a,w_a)$, the matrices $w_aw_b^\top+w_bw_a^\top$ are eigenvectors of $\mathcal L^{-1}\mathcal K$ with eigenvalues $\frac{\lambda_a\lambda_b}{\lambda_a+\lambda_b}\le\frac{\lambda_d}2$. Finally \eqref{eq:B4-fluct} reads $(\gamma\mathcal L-\gamma^2\mathcal K)\Sigma=\gamma^2\mathcal K(\Xi)$, because $\Sigma-\frac1n\sum_i(I-\gamma H_i)\Sigma(I-\gamma H_i)=\gamma\mathcal L(\Sigma)-\gamma^2\mathcal K(\Sigma)$: on an eigenvector of $\mathcal L^{-1}\mathcal K$ with eigenvalue $\mu$ the map is the multiplication by $\frac{\gamma\mu}{1-\gamma\mu}$, increasing in $\mu$ on $\gamma\mu<1$.
\end{proof}

\noindent \emph{Heuristic two-term law.} Replacing $g_{\max}$ by $\gamma\tilde L$ and $c_{\min}$ by $\gamma n\lmin(\bar H)$ in conclusion (2) of Section~\ref{sec:B4-comparison} gives, for a sum with unequal Hessians,
\begin{equation}\label{eq:B4-twoterm}
\begin{aligned}
\text{reshuffled:}\quad& r\approx\min\Big\{r^{\rm rr}\big(\gamma n\lmin(\bar H)\big),\ \log\frac{2-\gamma\tilde L}{\gamma\tilde L}\Big\},\\
\text{i.i.d., expectation, large }c:\quad& r\approx\frac{2(1-\gamma\tilde L)}{2-\gamma\tilde L},
\end{aligned}
\end{equation}
and the prediction that both orders lose mean-square stability at $\gamma\tilde L=1$. Only Proposition~\ref{prop:B4-Leff} is proved; \eqref{eq:B4-twoterm} assumes the slaved regime and treats the slow mode and the dispersion as decoupled, which is exact for identical Hessians and not otherwise (when the two terms are close the measured rate is below both). It is tested on logistic regression in the Outcome, where it accounts for the failure of the law $r^{\rm rr}(c_{\min})$ at large steps; $\tilde L$ is there well below $R^2$ ($0.16R^2$ to $0.43R^2$ over the twelve problems, i.e.\ $0.26$ to $2.98$ times the mean per-sample curvature at the optimum), which is why steps up to $\gamma R^2=1$ remain usable.
% B4b_summary.json: logistic.<problem>.{L_eff_over_R2, L_eff_over_mean_Li}

\paragraph{The orthogonal extreme.} The opposite of identical Hessians is a sum whose components act on orthogonal directions: $d=n$ and $f_i(\theta)=\frac{n\lambda_i}2(\theta_i-\theta^\ast_i)^2$, so that $F''=\diag(\lambda_i)$, $R^2=n\max_i\lambda_i$ and $\kappa=R^2/\mu\ge n$: this model lives entirely in the regime $\kappa\ge n$.

\begin{proposition}[orthogonal components]\label{prop:B4-orth}
For this sum, each coordinate $i$, with its own memory slot, evolves alone: with $e=\theta_i-\theta_i^\ast$, $\psi$ the error at the last visit of index $i$ and $c=\gamma n\lambda_i$, a step on another index gives $e\leftarrow e-\frac cn\psi$, and a step on index $i$ gives $e\leftarrow(1-c)e+c(1-\frac1n)\psi$, $\psi\leftarrow e$ (the error before the step). Consequently:
\begin{enumerate}[label=(\alph*)]
\item for a fixed order, one pass (ending with the visit) is the map $(e,\psi)\mapsto\begin{pmatrix}1-c&c^2\frac{n-1}n\\1&-c\frac{n-1}n\end{pmatrix}(e,\psi)$ (\texttt{orth\_cyclic\_matrix}), which is stable if and only if $c<\frac{2n}{3n-2}$, i.e.\ $c<\frac23$ in the limit; in the limit its spectral radius is smallest at $c=\frac12$, where the squared error contracts by $\log2=0.693$ nats per pass;
\item for the reshuffled order, in the limit $n\to\infty$ the visit takes place at a uniform time $u\in[0,1]$ of each pass and the \emph{mean} pass map is $\begin{pmatrix}1-\frac32c&\frac c2+\frac23c^2\\1&-\frac c2\end{pmatrix}$, stable if and only if $c<18-\sqrt{300}=0.6795$;
\item for i.i.d.\ indices a coordinate never sampled until time $t=un$ has $e=e_0(1-cu)$, a linear drift and no geometric contraction: such coordinates carry $e^{-u}(1-cu)^2e_0^2$ of the expected squared error, so the ceiling of one nat per pass holds for the expectation here too.
\end{enumerate}
\end{proposition}
\begin{proof}
$f_j'(\theta)$ has only the coordinate $j$, equal to $n\lambda_j(\theta_j-\theta_j^\ast)$, so the coordinate $i$ of the SAGA direction is $\bar z_i=\lambda_i\psi$ when $j\ne i$ and $n\lambda_i(e-\psi)+\lambda_i\psi$ when $j=i$; multiply by $\gamma$. (a) In a pass the $n-1$ other steps give $e_1=e-c\frac{n-1}n\psi$, then the visit gives $(1-c)e_1+c(1-\frac1n)\psi$ and $\psi\leftarrow e_1$; the matrix has trace $1-c-c\frac{n-1}n$ and determinant $-c\frac{n-1}n$, and its characteristic polynomial $p$ satisfies $p(1)=c>0$ and $p(-1)=2-c-2c\frac{n-1}n$: the roots (real, of opposite signs) are in $(-1,1)$ iff $p(-1)>0$. In the limit the roots are $\frac12(1-2c\pm\sqrt{1+4c^2})$, of equal modulus $\frac1{\sqrt2}$ at $c=\frac12$, which minimizes the larger modulus. (b) With the visit at time $u$: $e_1=e-uc\psi$, then $e'=(1-c)e_1+c\psi-(1-u)c\,e_1=(1-2c+uc)(e-uc\psi)+c\psi$ and $\psi'=e_1$; integrate over $u$. The matrix has trace $1-2c$ and determinant $-c+\frac{c^2}{12}$; $p(1)=c+\frac{c^2}{12}>0$ and $p(-1)=2-3c+\frac{c^2}{12}$ vanishes at $18-\sqrt{300}$. (c) As long as index $i$ is not sampled, $\psi=e_0$ and $e$ decreases by $\frac cne_0$ per step.
\end{proof}

\noindent In (b) the mean does not describe a run: each coordinate has a single slot, nothing averages, and the randomness of the visit time is an effect of order one. The typical exponents (Lyapunov exponents of the products of the random $2\times2$ maps, \texttt{orth\_exponent}, $n=1000$, $4\cdot10^4$ passes, three seeds) are:
% B4b_summary.json: orthogonal.by_c.<c>.{cyclic, rr.median, iid.median}
\begin{center}
\footnotesize
\begin{tabular}{@{}lccccccc@{}}
\toprule
$c=\gamma n\lambda=\gamma R^2$ & $0.05$ & $0.1$ & $0.25$ & $0.4$ & $0.5$ & $0.6$ & $0.7$ \\
\midrule
fixed order (exact, $n\to\infty$) & $0.097$ & $0.189$ & $0.424$ & $0.601$ & $0.693$ & $0.253$ & $-0.117$ \\
reshuffled, typical run & $0.098$ & $0.191$ & $0.436$ & $0.632$ & $0.694$ & $0.308$ & $-0.065$ \\
i.i.d., typical run & $0.100$ & $0.201$ & $0.496$ & $0.652$ & $0.585$ & $0.355$ & $0.023$ \\
\bottomrule
\end{tabular}
\end{center}
\noindent The three orders have exponents of the same order, close to the gradient rate $2c$ at small steps, and all lose stability near $\gamma R^2=\frac23$ to $0.7$: in this extreme the order hardly matters, the binding constraint is the per-sample curvature, and no exponent exceeds $0.7$ nats per pass. The two solvable extremes bracket the experiments: the order matters when the components are aligned ($\kappa\ll n$, the memory regime) and not when they are orthogonal ($\kappa\ge n$); and the stability threshold of a fixed order, $\gamma n\lmax(F'')<17.70$ in the first, becomes $\gamma R^2<\frac23$ in the second. The value $c_\ast$ is a property of the model, the existence of a threshold is not.
% B4b_summary.json: note.orth_cyclic_matrix_limit_c_two_thirds, note.orth_rr_mean_map_threshold

\subsection{Related work}\label{sec:B4-related}

\emph{Observations as old as the methods.} The two phenomena studied here were reported when these methods were introduced. For Finito, \cite[Section~3]{defazio2014finito} write that neither SAG nor Finito works ``when the same ordering is used each pass'', that pre-permuting the data and then making in-order passes does not work either, that reducing the step by one or two orders of magnitude does not fix it but that reducing it ``by a factor of $n$'' does, and that sampling without replacement within each pass gives no speed-up for SAG but ``speedups of up to a factor of two'' for Finito, with ``no theory to support this case''. For SAGA, \cite[Section~3]{defazio2014saga} report ``empirical speed-ups of up-to a factor of 2x'' for a per-pass re-permuted ordering, and \cite[Section~5.2]{defazio2014new} that the cyclic ordering ``converges but extremely slowly'' while the permuted ordering is sometimes $50\%$ slower or faster than the randomized one. For the incremental aggregated gradient method \cite{blatt2007convergent}, \cite{schmidt2017minimizing} note that a step $\frac1{16L}$ may diverge and that a step of about $\frac1{nL}$ is needed, the admissible step being ``not precisely characterized''. The statements (R2) and (R3) and the fixed-step laws are exact counterparts, on identical-Hessian sums, of these observations: a fixed order needs $\gamma n\lmax<c_\ast(n)$, i.e.\ a step $n$ times smaller; reshuffling gains a factor up to $2$ over the \emph{expectation} of the i.i.d.\ order and a smaller, step-dependent factor over its typical runs (Section~\ref{sec:B4-typical}).

\emph{Rates proved without replacement.} \cite{ying2017variance} give the first proof of linear convergence of SAGA under random reshuffling (their Theorem~1; they note that no proof existed for variance-reduced methods under reshuffling): in our notation, for $\gamma\le\frac{\mu}{11R^4n}$ a Lyapunov function contracts by $(1-\frac{\gamma\mu n}4)/(1-\frac{27R^8\gamma^3n^3}{\mu})$ per pass, i.e.\ about $1-\frac1{410\kappa^2}$ at the largest admissible step and $1-\frac1{108\kappa^2}$ at $\gamma=\frac\mu{18R^4n}$, a number of passes of order $\kappa^2\log\frac1\varepsilon$ which they describe as slower than the rate proved for SAGA in \cite{defazio2014saga}; they also state the mechanism studied here, in words (reshuffling ``updates every history variable during each run, while uniform sampling may leave some variables [\dots] untouched''), and observe that reshuffled SAGA is faster in their simulations. \cite{park2020linear} prove linear convergence of cyclic SAGA when every $f_i$ is $\mu$-strongly convex and $R^2$-smooth, for $\gamma\le\frac{\mu}{65\sqrt{n(n+1)}R^4}$, with the factor $1-\frac1{368\kappa^2}$ per pass at half this step, and report that in their experiments random permutations outperform i.i.d.\ sampling, which outperforms the cyclic order, the study of this effect being left to future work. For Finito-type methods the rates are those of gradient descent: the cyclic method DIAG of \cite{mokhtari2018surpassing} contracts $\norm{\theta-\etastar}$ by $\frac{\kappa-1}{\kappa+1}$ per pass with the step $\frac2{\mu+L}$ (their Corollary~3; every $f_i$ being $\mu$-strongly convex and $L$-smooth, $\kappa=L/\mu$), and \cite{huang2021improved} obtain for a damped Finito, under reshuffling and cyclic sampling and the same assumption on every component, the factor $1-\frac{2\theta\alpha\mu L}{\mu+L}$ per pass on the squared error for a step $\alpha\le\frac2{\mu+L}$ and a damping $\theta\in(0,1)$ (their Theorem~3), i.e.\ $(\frac{\kappa-1}{\kappa+1})^2$ at best; they state that the existing rates for variance reduction without replacement, theirs included, do not match those of i.i.d.\ sampling (except, in the convex case, for an optimal cyclic order in a highly heterogeneous scenario), closing this gap being left open. \cite{malinovsky2023random} analyse SVRG under reshuffling, shuffling once and the cyclic order, each epoch using a full gradient: epoch complexities of order $\kappa^{3/2}$, of order $\kappa$ when $n\gtrsim\kappa$, and $\kappa\sqrt{\kappa/n}$ for strongly convex components and $n$ large (arXiv abstract and Theorems~1--3), in the line of the analysis of reshuffled SGD of \cite{mishchenko2020random}. For the incremental aggregated gradient method, \cite{gurbuzbalaban2017convergence} prove a linear rate, and \cite{vanli2016stronger} improve it to $1-O(\frac1{\kappa n})$ per iteration with steps of order $\frac1{nL}$: in the units of this note a sufficient $c$ of order $\frac13$, a constant factor below the model threshold $\frac{\pi^2}2$ of Remark~\ref{rem:B4-visible}(d). In order of magnitude all these rates are weaker than the i.i.d.\ ones ($n\kappa$ gradients at best against $n+\kappa$). With their constants one of them is not: $(\frac{\kappa-1}{\kappa+1})^2<e^{-1}$ for $\kappa<\frac{\sqrt e+1}{\sqrt e-1}=4.08$, so for such $\kappa$ cyclic or reshuffled Finito with one gradient per iteration is \emph{proved} to contract by more than one nat per pass, hence to go below the floor of Proposition~\ref{prop:B4-floors}(a) for span methods with i.i.d.\ indices, a comparison that these papers do not make.

\emph{Coordinate descent.} For coordinate descent with exact line search on quadratics, the setting dual to these table methods, the comparison of the three orders is a theorem: \cite{gurbuzbalaban2020randomness} compute the exact worst-case rates under random permutations, i.i.d.\ sampling and a fixed order on quadratics with diagonally dominant Hessians and prove that a fixed order is faster than random permutations, which are faster than i.i.d.\ sampling, with a gain that grows as the Hessian becomes more diagonally dominant; \cite{lee2019random} give a tight analysis of random permutations on a worst case of the cyclic order. Their device, an exact expected pass matrix followed by Jensen's inequality, is that of Proposition~\ref{prop:B4-rr}. What differs here is the memory of gradients: it is what makes a fixed order unstable (the opposite ranking) and what produces the ceilings of one and two nats.

\emph{Lower bounds.} \cite{hannah2018breaking} show that the term $n\log\frac1\varepsilon$ of the bound of \cite{lan2018optimal} rests on the span assumption with i.i.d.\ indices and that SVRG with full gradients and a tuned epoch length and step size goes below it by the factor $1+(\log\frac n\kappa)_+$, which their lower bound for a larger class shows to be the right order (Theorem~\ref{thm:HLOY}).

\emph{What was not found} (searches of November 2026 on reshuffled, cyclic and shuffled-once SAGA, Finito and incremental aggregated gradient methods, on permutation orders in coordinate descent and on delay equations; the reading depth of each work is given in Section~10 of the literature review). (a) An instability threshold of fixed-order SAGA such as $c_\ast$: not found; the conditions quoted above for SAGA are sufficient ones, with steps of order $\frac1{nR^2\kappa}$. For the incremental aggregated gradient method the threshold $\frac{\pi^2}2$ is the classical stability boundary of the limiting delay equation \cite{matsunaga2011asymptotic}; we did not find it stated for the method. (b) Exact rates of SAGA under reshuffling, for the mean or for the mean square: not found; the closest results are the exact rates of coordinate descent under permutations and the worst-case factor of cyclic Finito. (c) A result in which reshuffling provably beats i.i.d.\ sampling for a variance-reduced method: this holds for coordinate descent on quadratics \cite{gurbuzbalaban2020randomness}, and for Finito with $\kappa<4.08$ by the combination made above; it was not found for SAGA, nor for larger $\kappa$. (d) A method with one gradient per iteration and no replacement going below a lower bound proved for i.i.d.\ sampling: the same corner $\kappa<4.08$ is the closest result; \cite{hannah2018breaking} do it with full gradients. The present note does not fill these gaps in general: its statements are exact for identical Hessians only, where they give SAGA's rates for every $\kappa$ and both orders. The fixed-step laws, the effective curvature and the heuristic for typical runs were not the object of a dedicated search.

\subsection{Limits and outcome}\label{sec:B4-outcome}

\paragraph{Limits.} (i) Everything proved rests on Assumption~\ref{ass:B4-model}: with different Hessians the modes are coupled; the two-term law \eqref{eq:B4-twoterm} is a heuristic, tested on logistic regression. Under this assumption the component gradients differ by constants, so a control variate taken at one common reference point is exact: the directions of SVRG and L-SVRG equal the full gradient and their error decays at $2c$ nats per pass with no memory ceiling. The ceilings of one and two nats are those of a \emph{table} of gradients taken at different points, and the model cannot compare the orders of SAGA with reference-point methods. (ii) The statements on rates are limits $n\to\infty$, at fixed $c$ or at fixed $g$ (Section~\ref{sec:B4-model}); the corrections are $O(\frac1n)$ for (R1) and (R3) and large for the critical step of (R2) ($K/n$ with $K\approx352$). (iii) For the reshuffled order the fixed-$g$ limit is derived and verified on the exact matrix, not proved; that typical runs follow the mean-square rate is numerical. (iv) The model does not contain the instance of Section~\ref{sec:B4-lower}. (v) The typical rate under i.i.d.\ sampling is measured; \eqref{eq:B4-typical} is a heuristic and its limit a conjecture.

\paragraph{Outcome of the experiments.}
% Every number of this paragraph is read in results/B4_summary.json (experiments/exp_B4_sampling_order.py, Figure B4)
% or in results/B4b_summary.json (experiments/exp_B4b_refinements.py, Figure B4b); their entries "definitions" define
% each quantity. Details and limits: results/findings.md, Section 9.
The claim holds in the form ``the ceiling of one nat per pass concerns the \emph{expectation} under sampling with replacement'', and the gain of reshuffling over what an i.i.d.\ run actually does is a factor $1.2$ to $1.5$ in the exponent on logistic regression with $\kappa\le n$, not $2$.

% B4: model.by_c.<c>.n=<n>.{rr_simulated, rr_mean_square, rr_first_moment, iid_expectation, iid_typical, rr_over_iid_typical};
% model.n=<n>.*; model.expectation_vs_runs
\emph{Model} ($n\in\{100,10^3,10^4\}$, $c\in\{\frac14,\dots,32\}$, $9$ seeds, exponents fitted on passes $30$--$150$). \emph{Reshuffling:} the simulated exponents (medians over the seeds) are $1.1881$, $1.5380$, $1.7534$, $1.9243$ at $c=1,2,4,16$ for $n=10^3$, against $r^{\rm rr}=1.1874$, $1.5370$, $1.7523$, $1.9230$; over the eight values of $c$ the largest deviation from the exact mean-square exponent $r^{\rm rr}_n(c)$ of Proposition~\ref{prop:B4-rr-ms} is $0.016$ ($n=100$), $3\cdot10^{-4}$ ($n=10^3$) and $7\cdot10^{-5}$ ($n=10^4$). The exact matrix also accounts for the one cell where the fluid law fails: for $n=100$, $c=32$ ($g=0.32$) the simulated exponent is $1.619$ and $r^{\rm rr}_n=1.603$, whereas the squared mean and $r^{\rm rr}$ give $1.969$ and $1.959$. Typical reshuffled runs thus decay at the mean-square rate. \emph{I.i.d.\ indices:} the exact exponent of $\E e_t^2$ is $0.9995$, $0.9992$, $0.9982$, $0.9922$ at $c=1,2,4,16$ ($n=10^3$), but single runs decay faster (Section~\ref{sec:B4-typical}): for $n=10^3$, $c=2$, the exact $\E e_t^2$ after $20$ passes is $3.7\cdot10^{-12}$, the median of the nine runs $7.3\cdot10^{-17}$ and their mean $8.0\cdot10^{-16}$ (the last initial slot is refreshed after $7.7$ passes in median): a seed average does not estimate the expectation. \emph{Consequence:} at equal $c$ and $n=10^3$ the reshuffled exponent is $0.74$ and $0.97$ times the typical i.i.d.\ one at $c=1,2$, and $1.16$, $1.27$, $1.35$, $1.46$ times at $c=4,8,16,32$. \emph{Fixed order:} the companion matrix gives $c_\ast(100)=14.784$ and $c_\ast(1000)=17.36$, against $c_\ast=17.702$.

% B4: logistic.headline.<problem>.{iid_textbook, iid_best, rr_best, gain}; logistic.law_R3.*; logistic.fixed_orders.*;
% logistic.status_counts; B4b: logistic.<problem>.by_step, logistic.summary, logistic.<problem>.L_eff_over_R2
\emph{Logistic regression} ($n=1000$, $d=50$, $\kappa\in\{10,10^2,10^3,10^4\}$, $\alpha\in\{0,1,2\}$, one data set per problem, $5$ sampling seeds, steps $\gamma R^2\in\{\frac1{128},\frac1{64},\dots,2\}$, $300$ passes; exponents of $\norm{\theta_t-\etastar}^2$ fitted on the second half of each trajectory above $10^{-18}$, medians over the seeds, ranges over $\alpha$; ``best'' is the best step of the grid for each order; passes to $10^{-10}$ at the fastest step of each order, i.i.d.\ $\to$ reshuffled, for the first passage below the target and for the pass from which the error stays below it).
\begin{center}
\footnotesize
\begin{tabular}{@{}lcccc@{}}
\toprule
$\kappa/n$ & i.i.d., $\gamma=\frac1{4R^2}$ & i.i.d., best step & reshuffled, best step & ratio of bests \\
\midrule
$0.01$ & $1.26$--$1.30$ & $1.50$--$1.61$ & $1.99$--$2.01$ & $1.23$--$1.33$ \\
$0.1$ & $1.13$--$1.30$ & $1.34$--$1.39$ & $1.88$--$1.93$ & $1.35$--$1.41$ \\
$1$ & $0.53$--$1.06$ & $1.06$--$1.07$ & $1.44$--$1.54$ & $1.36$--$1.46$ \\
$10$ & $0.065$--$0.49$ & $0.48$--$0.93$ & $0.43$--$1.01$ & $0.90$--$1.08$ \\
\bottomrule
\end{tabular}

\smallskip
\begin{tabular}{@{}lcccc@{}}
\toprule
$\kappa/n$ & first passage below $10^{-10}$ & saving & stays below $10^{-10}$ from & saving \\
\midrule
$0.01$ & $11.6$--$12.7\to10.1$ & $13$--$20\%$ & $12.5$--$13.4\to12.0$ & $4$--$11\%$ \\
$0.1$ & $13.9$--$16.3\to11.1$--$12.1$ & $20$--$26\%$ & $13.9$--$16.3\to12.0$--$12.1$ & $13$--$26\%$ \\
$1$ & $18.9$--$19.7\to14.2$--$16.2$ & $18$--$26\%$ & $19.3$--$19.7\to14.2$--$16.2$ & $18$--$26\%$ \\
$10$ & $23.1$--$40.6\to21.3$--$45.3$ & $-12$--$8\%$ & $23.1$--$40.6\to21.3$--$45.3$ & $-12$--$8\%$ \\
\bottomrule
\end{tabular}
\end{center}
\noindent For $\kappa\le n$ reshuffling at its best step raises the exponent by $1.23$--$1.46$ over i.i.d.\ sampling at its own best step; for $\kappa=10n$ there is no gain (ratios $1.08$, $0.94$, $0.90$ for $\alpha=0,1,2$). The saving in passes depends on how the passage is counted, because under reshuffling the error is not monotone within a pass (it dips shortly after each pass boundary and rises until the next one; the same sawtooth is seen in the model and on the hard instance): $13$--$26\%$ of the passes to the \emph{first} passage below $10^{-10}$, but $4$--$26\%$ of the passes from which the error \emph{stays} below it ($4$--$11\%$ at $\kappa/n=0.01$, where the first passage at $10.1$ passes is a dip and the error stays below the target from $12.0$ passes on). The i.i.d.\ runs themselves exceed one nat per pass ($1.13$--$1.30$ at the textbook step for $\kappa\le n/10$), as the typical runs of the model do. \emph{Law (R3)}, evaluated at $c=\gamma n\lmin(F''(\etastar))$ as a heuristic: on the $65$ (problem, step) pairs with $\gamma R^2\le\frac12$ at which every reshuffled run converged, the measured exponent is within $8.6\%$ of the law (median deviation $+2.2\%$); on the $24$ pairs with $\gamma R^2\ge1$ it is within $10\%$ for $11$ pairs only and falls $91\%$ below the law in the worst case ($\kappa=10^4$, $\alpha=0$, $\gamma R^2=2$). \emph{Two-term law} \eqref{eq:B4-twoterm}: the dispersion branch $\log\frac{2-\gamma\tilde L}{\gamma\tilde L}$ accounts for these failures. On the $24$ pairs with $\gamma R^2\ge1$ the median error falls from $13\%$ to $3.8\%$ and $19$ pairs are within $10\%$; on the $8$ of them with $\kappa\le100$ every pair is within $6.6\%$ (the law alone: up to $54\%$). For $\kappa\ge10^3$ the two terms are close and the measured exponent is below both (errors up to $56\%$): the law is an upper envelope there, not a prediction. Over the $89$ pairs of the grid, $84$ are within $10\%$ of the two-term law, against $76$ for the law (R3) alone. The step at which the two terms cross, $\gamma R^2$ between $0.8$ and $2$ for $\kappa\le10^3$, is where the best step of the grid sits ($\frac12$ to $2$). \emph{Fixed orders:} the shuffled-once and cyclic orders fail (divergence, or $10^{-10}$ not reached in $300$ passes with an exponent below half the i.i.d.\ one at the same step) on the same $43$ of the $108$ (problem, step) pairs: the smallest $\gamma n\lmax(F''(\etastar))$ with a failure is $8.0$, the largest without one $10.9$, no pair above $c_\ast$ is free of failure and $8$ failing pairs lie below it ($c_\ast(1000)=17.36$); counting also the pairs that converge within the budget with an exponent below half the i.i.d.\ one, $55$ pairs are slow or failing, from $\gamma n\lmax=4.3$ on, $20$ of them below $c_\ast$. These are the slow-downs that the model predicts before the instability (fixed-order exponent $0.19$ at $c=8$ and $0.057$ at $c=12$ for $n=1000$). No fixed-order run exploded: the $91$ shuffled-once and $19$ cyclic runs classified as diverged ended above their initial error after $300$ passes without reaching $10^{12}$ times it, and $92$ and $17$ further runs were still growing, which is the weak instability of Remark~\ref{rem:B4-visible}(a); no i.i.d.\ or reshuffled run grew. The exponent of (R2), taken at the worse end of the spectrum of $F''(\etastar)$, differs from the measured fixed-order exponents by $0.008$ nats per pass in median and by at most $0.21$--$0.22$.

% B4: hard_instance.kappa/n=<r>.by_eps.<eps>.{bound_iid, bound_any, iid.best, rr.best, caveat_iid}; hard_instance.kappa/n=<r>.{q, below_iid_floor}
\emph{Hard instance} ($n=256$, $m=12$, $32$ seeds, steps $\gamma R^2\in\{\frac18,\dots,4\}$; gradient count at which the mean over the seeds of the relative squared error reaches $\varepsilon$, best step of the grid, divided by the bound of Corollary~\ref{cor:sweep}; for the reshuffled order, whose mean is a sawtooth, three conventions: first passage / count from which the mean stays below $\varepsilon$ / first pass end below $\varepsilon$).
\begin{center}
\footnotesize
\begin{tabular}{@{}lcccccc@{}}
\toprule
 & \multicolumn{3}{c}{reshuffled SAGA (first / stays / pass end)} & \multicolumn{3}{c}{SAGA with i.i.d.\ indices} \\
$\kappa/n$ & $\varepsilon=10^{-8}$ & $10^{-16}$ & $10^{-24}$ & $\varepsilon=10^{-8}$ & $10^{-16}$ & $10^{-24}$ \\
\midrule
$0.1$ & $0.90$ / $0.99$ / $1.08$ & $0.65$ / $0.68$ / $0.71$ & $0.61$ / $0.63$ / $0.65$ & $1.20$ & $0.89$ & $0.82$ \\
$\frac13$ & $0.90$ / $0.90$ / $0.97$ & $0.74$ / $0.76$ / $0.78$ & $0.72$ / $0.76$ / $0.74$ & $1.11$ & $0.95$ & $0.87$ \\
$1$ & $1.37$ / $1.37$ / $1.41$ & $1.26$ / $1.26$ / $1.28$ & n/a & $1.50$ & $1.40$ & n/a \\
$4$ & $3.22$ / $3.22$ / $3.24$ & n/a & n/a & $3.03$ & n/a & n/a \\
\bottomrule
\end{tabular}
\end{center}
\noindent (n/a: $m=12$ does not satisfy the dimension condition $q^{2m}\le\varepsilon$.) The statement that does not depend on a convention is a comparison at equal counts: for $\kappa/n=0.1$, from $10.03\,n$ gradients on, and for $\kappa/n=\frac13$ from $9.06\,n$ on, the seed mean of reshuffled SAGA is below the floor \eqref{ineq:expected-sweep} that the \emph{expected} error of every span method with i.i.d.\ indices must exceed at the same count, at every record down to the rounding level. In counts, for $\kappa\le n/3$ reshuffled SAGA needs $0.61$--$0.78$ of the bound at $\varepsilon=10^{-16}$ and $10^{-24}$ whatever the convention; at $10^{-8}$ it first passes below $\varepsilon$ at $0.90$ of the bound ($10.06\,n$ against $11.14\,n$ and $12.08\,n$ against $13.45\,n$) but, for $\kappa/n=0.1$, stays below it only from $0.99$ of the bound and is below it at a pass end at $1.08$. The best reshuffled step is interior to the grid ($\gamma R^2=2$ and $1$). These counts stay $3.4$ to $4.3$ times above the bound of Proposition~\ref{prop:anyseq} for $\kappa\le n/3$ ($6.8$ at $\kappa=4n$), and reshuffling does not go below the i.i.d.\ bound for $\kappa\ge n$ (at $\varepsilon=10^{-8}$ it needs $0.91$ times the gradients of i.i.d.\ SAGA for $\kappa=n$ and $1.06$ times for $\kappa=4n$). \emph{Seed variability:} over $500$ bootstrap resamples of the seeds the reshuffled counts move by at most $0.02\,n$ for $\kappa/n=0.1$, whereas the i.i.d.\ count at $10^{-8}$ ranges over $12.25$--$13.76\,n$ ($1.10$--$1.23$ times the bound), so that reshuffled SAGA needs $0.73$--$0.82$ times the gradients of i.i.d.\ SAGA there. \emph{Caution:} under i.i.d.\ sampling the seed mean of SAGA itself falls below the bound at $10^{-16}$ and $10^{-24}$, which contradicts nothing: at the corresponding counts the expectation of the pathwise floor \eqref{ineq:tail}, a certified lower bound on the expected error, is $47$ and $9.7\cdot10^3$ times the seed mean for $\kappa/n=0.1$ ($9.4$ and $1.3\cdot10^3$ times for $\kappa/n=\frac13$), one seed carries $0.91$--$0.92$ of the mean ($0.50$ and $0.72$), and no (seed, component) pair was left unqueried, against $7\cdot10^{-8}$ and $1.5\cdot10^{-13}$ expected for $\kappa/n=0.1$. At $\varepsilon=10^{-8}$ the certified floor is $0.23$ and $0.46$ times the seed mean, so the i.i.d.\ column is not below it there; for reshuffled SAGA, whose floor is pathwise, the largest share of one seed in the mean is $0.07$--$0.16$, except $0.52$ at $\kappa/n=\frac13$, $\varepsilon=10^{-24}$. One $n$, one block length and a grid of six steps limit these statements.
